\documentclass[11pt]{article}
\usepackage[T1]{fontenc}
\usepackage{lmodern,amsmath,amssymb,amsthm,mathtools,booktabs,microtype}
\usepackage[margin=1.02in]{geometry}
\usepackage[hidelinks,unicode]{hyperref}
\usepackage{xurl}
\hypersetup{
 pdftitle={Subform Zeta Functions of Integral Quadratic Lattices in Arbitrary Rank},
 pdfauthor={Zhanhe Zhang},
 pdfsubject={Intrinsic isometry classes and zeta functions of integral quadratic lattices},
 pdfkeywords={quadratic lattices, subform zeta functions, spinor genera, orientation}
}
\newtheorem{theorem}{Theorem}[section]
\newtheorem*{mainthmA}{Theorem A (positive definite lattices)}
\newtheorem*{mainthmB}{Theorem B (orientation and parity)}
\newtheorem*{mainthmC}{Theorem C (indefinite lattices)}
\newtheorem{lemma}[theorem]{Lemma}
\newtheorem{corollary}[theorem]{Corollary}
\newtheorem{proposition}[theorem]{Proposition}
\theoremstyle{remark}\newtheorem{remark}[theorem]{Remark}
\setlength{\emergencystretch}{2em}
\title{Subform Zeta Functions of Integral Quadratic Lattices\\in Arbitrary Rank}
\author{Zhanhe Zhang}
\date{}
\begin{document}
\maketitle
\begin{abstract}
For a fixed nondegenerate integral quadratic lattice, we count
intrinsic isometry classes of finite-index sublattices under ordinary
and proper isometry. For positive definite lattices of rank $n\ge2$,
we obtain two exact coefficient formulas, one based on restoring
lattices at primes with cyclic discriminant group and the other on
signed incidence among rationally isometric parents. The resulting Dirichlet series extend meromorphically to $\Re s>n-1$
and have a unique simple pole there, at $s=n$. Their summatory coefficients
have main terms of order $X^n$, with errors
$O_{L,\varepsilon}(X^{n-1+\varepsilon})$.
For a nondegenerate lattice of arbitrary signature, ordinary and proper
counts agree at every index if and only if the rank is odd. In indefinite
rank at least three, local embedding spinor norms determine the represented
spinor classes. The resulting counting functions have
a common leading term of order $X(\log X)^2$ and a strictly positive
secondary term of order $X(\log X)^{3/2}$; neither series has a
meromorphic germ at $s=1$. For the cubic ternary lattice, an all-index formula exhibits a
square-root singularity at $s=2$.
\end{abstract}
\noindent\textit{2020 Mathematics Subject Classification.} Primary 11E12, 11E16, 11M41; Secondary 11H06.\\
\noindent\textit{Keywords.} Integral quadratic lattices, subform zeta functions, intrinsic isometry, lattice embeddings, spinor genera, orientation.

\section{Introduction}

\subsection{Intrinsic isometry classes and earlier work}

Let $L$ be an integral lattice in a nondegenerate quadratic space $(V,Q)$
over $\mathbb Q$. We count the isometry classes represented by the
restrictions of $Q$ to finite-index subgroups $M\subseteq L$. Such
isometries need not preserve $L$, so their classes may identify distinct
orbits under the automorphism group of the parent.

Fix an orientation of $V$ and let $H$ denote either $O(V)(\mathbb Q)$
or $SO(V)(\mathbb Q)$. Write $a_H(L;m)$ for the number of represented
$H$-isometry classes at index $m$, and set
\begin{equation}\label{eq:introdef}
 Z_H(L;s)=\sum_{m\ge1}\frac{a_H(L;m)}{m^s},\qquad
 \zeta_L^{\rm GL}(s)=Z_O(L;s),\qquad
 \zeta_L^{\rm SL}(s)=Z_{SO}(L;s).
\end{equation}
We suppress $L$ when fixed. Ordinary and proper isometries have
determinant $\pm1$ and $1$, respectively; a represented class has a
representative contained in $L$. In rank one both series equal
$\zeta(s)$. Negating a definite form does not change either count.

For $L=\mathbb Z^3$ with its standard quadratic form, the proper
automorphism group has $241$ orbits on index-$71$ sublattices,
representing only $240$ intrinsic classes (Theorem~\ref{thm:cubic-all}).
The remaining identification comes from a rational isometry moving $L$.

Kim, Lee and Lee introduced these series for positive definite
binary forms of fundamental discriminant \cite{KLL} and subsequently
studied the rational hyperbolic plane \cite{Hyperbolic}.
The nonmaximal binary case under a generic-unit hypothesis was treated
by Zhang \cite{Zhang}.
\subsection{Main results}

\begin{mainthmA}
Let $L$ be a positive definite integral lattice of rank $n\ge2$ and
$H\in\{O(V)(\mathbb Q),SO(V)(\mathbb Q)\}$. For every $m\ge1$,
$a_H(L;m)$ is given by both the intrinsic core-restoration formula of
Theorem~\ref{th:core} and the finite parent-incidence formula of
Theorem~\ref{th:incidence}. Moreover,
\[
 Z_H(L;s)=\zeta(s-n+1)D_H(s)+E_H(s),
 \qquad \Re s>n-1,
\]
where $D_H$ and $E_H$ are represented by absolutely convergent Dirichlet
series in this half-plane. The abscissa of convergence of $Z_H$ is $n$;
its only pole in $\Re s>n-1$ is simple, at $s=n$, with residue $R_H>0$.
For every $\varepsilon>0$,
\begin{equation}\label{eq:intropositive}
 \sum_{m\le X}a_H(L;m)
 =\frac{R_H}{n}X^n+O_{L,\varepsilon}(X^{n-1+\varepsilon}).
\end{equation}
\end{mainthmA}

For $n\ge3$, the fixed-parent coefficient sequences are
nonmultiplicative (Theorem~\ref{th:nonmult}), although the scalar
genus mass admits an Euler product (Section~\ref{sec:genusmass}).

In odd rank, $-I$ reverses orientation while preserving every lattice.
In even rank, a chiral ordinary class may have two proper orientations
with only one represented in $L$.

\begin{mainthmB}
Let $L$ be any nondegenerate integral quadratic lattice of positive rank
$n$, of arbitrary signature. Then
\[
 a_O(L;m)=a_{SO}(L;m)\text{ for every }m\ge1
 \quad\Longleftrightarrow\quad n\text{ is odd}.
\]
If $n$ is even, there exist a positive integer $b$ and a finite set of
primes $S'$ depending on $L$ such that
\[
 a_{SO}(L;bt)>a_O(L;bt)
 \qquad\text{whenever }\gcd\!\left(t,\prod_{p\in S'}p\right)=1.
\]
In particular, the two sequences differ on a set of indices of positive
natural density.
\end{mainthmB}

In even positive definite rank the residues satisfy
$\tfrac12R_{SO}\le R_O<R_{SO}$, with equality on the left precisely when
the parent admits an improper automorphism
(Theorem~\ref{th:orientation}). Section~\ref{sec:quaternary} gives
explicit even-rank families, including cases where orientation depends
on the choices at several primes.

In indefinite rank at least three, lattice stabilizers may be
infinite. Eichler's theorem gives one proper class per proper spinor
genus \cite[\S9.2]{SP}; local embedding spinor norms determine which
classes are represented in $L$.

\begin{mainthmC}
Let $L$ be an indefinite integral lattice of rank $n\ge3$, with either
choice $H=O$ or $SO$. Its represented-class coefficients at every index
are determined by local lattice types and the spinor-norm images of
admissible embeddings, as described in Section~\ref{sec:indef}. There are
constants $C_L>0$ and $K_H>0$ such that
\begin{equation}\label{eq:introindefinite}
 \sum_{m\le X}a_H(L;m)
 =\frac{C_L}{2}X(\log X)^2
 +\frac{K_H}{\Gamma(5/2)}X(\log X)^{3/2}
 +o_L\bigl(X(\log X)^{3/2}\bigr).
\end{equation}
The leading constant $C_L$ is the same for the two orientation
conventions, and neither corresponding Dirichlet series has a meromorphic
germ at $s=1$.
\end{mainthmC}

For squarefree $t$ coprime to $2\det B_L$, one has
$a_O(t)=a_{SO}(t)=3^{\omega(t)}$. In even indefinite rank at least four,
$K_{SO}>K_O$ (Theorem~\ref{th:indorient}).

\subsection{Methods and relation to other counting problems}

At good primes, a cyclic child discriminant group uniquely determines
an integral overparent. Restoration commutes with rational isometries,
reducing the count of regular extensions to the automorphism group of
the restored core (Section~\ref{sec:core}).

The signed incidence formula accounts for children contained in
several rationally isometric copies of $L$
(Theorem~\ref{th:incidence}). For the cubic lattice, local neighbor
trees and quaternionic rotations give explicit coefficients
\cite{Baake,BaakeZeiner,Chenevier,Zeiner,Voight}.

For indefinite lattices, character expansions of the represented
spinor-norm supports produce the $5/2$ exponent in Theorem C. The
definite local computations also use the theory of invariant-submodule
zeta functions \cite{Solomon,Rossmann}. The related reconstruction
problem for integral forms was studied by Chan and Oh \cite{ChanOh}.

Remark~1.7 and Conjecture~7.2 of the original Kim--Lee--Lee
preprint \cite{KLLpre} concern higher-rank extensions and analytic
properties, respectively. Corollary~\ref{cor:KLL} establishes the
scalar and nonscalar assertions. For the determinant-minus-one sum,
the conclusion depends on parity: it is holomorphic on $\Re s>n-1$
in even rank, whereas in odd rank the improper scalar $-I$ contributes
and the sum equals $Z_O/2$.

\subsection{Examples and organization}

The positive definite theory occupies Sections~\ref{sec:core}--\ref{sec:genusmass},
with ternary and even-rank applications in Sections~\ref{sec:cubic}
and~\ref{sec:quaternary}. Sections~\ref{sec:indef}--\ref{sec:realbinary}
treat indefinite lattices, followed by the parity criterion in
Section~\ref{sec:allparity}.

\section{Restoration at cyclic discriminant primes}\label{sec:core}
Let $V$ be a positive definite rational quadratic space of dimension
$n\ge2$, $L$ a full lattice, and $B$ an integral symmetric polarization obtained by scaling
the form. After fixing an orientation of $V$, write
$H=O(V)(\mathbb Q)$ or $SO(V)(\mathbb Q)$,
$\Gamma_H(U)=\{g\in H:gU=U\}$, and $c_H=|H\cap\{1,-1\}|$.
Write $a_H(m)$ for the corresponding represented-class count.

An isometry between full lattices extends uniquely over $\mathbb Q$;
it is proper precisely when its determinant in oriented bases is $+1$.
Positive rational rescaling of $B$ leaves the isometry groups and
indices unchanged. An index-$m$ restriction has Gram determinant
$m^2\det B_L$. The number $b_n(m)$ of raw index-$m$ subgroups
is given by Hermite normal form, with local series
$\prod_{j=0}^{n-1}(1-p^jT)^{-1}$.

For any finite set $S$ containing the prime divisors of $2\det B_L$, put $M^\#=\{x:B(x,M)\subset\mathbb Z\}$. Call $p^e\Vert [L:M]$ \emph{regular} if $p\notin S$ and $(M^\#/M)_p$ is cyclic, where $e>0$. Write $[L:M]=dt$ with $t$ the product of all the regular prime powers. Call $N$ a core if $N\subset L$ has no regular prime divisor of its index. Let $\mathcal C_H(d)$ consist of the $H$-classes of cores of index $d$ actually represented in $L$. Both definitions refer to the chosen $S$.

\begin{lemma}\label{lem:local}
At $p\notin S$, with $P=L_p$ unimodular, regular index-$p^e$ sublattices are the kernels of primitive covectors on $P/p^eP$ up to unit multiples whose reductions have nonisotropic normal. Each has a unique integral overlattice of relative index $p^e$ in $V_p$.
\end{lemma}
\begin{proof}
Let $r$ be the dimension of the image of $M_p$ in $P/pP$. The rank of the reduced Gram matrix of $M_p$ is at most $r\le n-1$. Cyclicity of the discriminant group, of order $p^{2e}$, is equivalent by Smith form to Gram rank $n-1$. This forces $r=n-1$, equivalently $P/M_p$ is cyclic, and the image hyperplane has nondegenerate restriction. Conversely these two conditions imply Gram rank $n-1$. The perfect pairing identifies the annihilator with a normal line, giving the covector description. Every integral overlattice is contained in $M_p^\#$, and a cyclic group of order $p^{2e}$ has exactly one subgroup of order $p^e$. Existence is supplied by $P$.
\end{proof}

When $p\nmid d$ and $p\notin S$, one has $N_p=L_p$. For
$h\in\Gamma_H(N)$, let $g_{p,h}(e)$ count the kernels in
Lemma~\ref{lem:local} fixed by $h$, with $g_{p,h}(0)=1$.
Equivalently, count primitive rows $v$ modulo $p^e$, up to unit
rescaling, for which $v B_{N}^{-1}v^t$ is a unit and
$vh\in(\mathbb Z/p^e)^\times v$. The inverse Gram matrix is taken
over $\mathbb Z_p$.
\begin{theorem}\label{th:core}
For all $m\ge1$ one has
\[
 a_H(m)=
 \sum_{\substack{dt=m,\ (d,t)=1\\p\mid t\Rightarrow p\notin S}}
 \sum_{[N]\in\mathcal C_H(d)}
 \frac1{|\Gamma_H(N)|}\sum_{h\in\Gamma_H(N)}
       \prod_{p^e\Vert t}g_{p,h}(e).
\]
\end{theorem}
\begin{proof}
Given $M$, enlarge $M_p$ at each regular prime by the unique subgroup in the lemma, leaving every other completion untouched. The unique global lattice with those completions is their intersection in $V$, equivalently the inverse image of the sum of the indicated primary subgroups in $M^\#/M$. The restored lattice lies in $L$ and is a core of index $d$.
The index, discriminant module and selected primary subgroups are
intrinsic to $M$, so restoration commutes with rational isometries.

Conversely, choose a representative $N\subset L$ of a core class. At primes dividing $t$ insert any of the indicated local kernels in $N_p$ and leave the other completions fixed. The Chinese remainder theorem gives a unique child in $N$, and restoration returns $N$. Every represented child with this core class is isometric under $H$ to one obtained from this representative, since transporting to $N$ places its image inside $N\subset L$. Two children of this same representative are equivalent exactly under $\Gamma_H(N)$. Burnside's lemma applied to the resulting product of local kernel sets gives the formula.
\end{proof}

The core-class sum also has a raw-lattice form. Let
$e_H(N,L)$ count integral overlattices $W$ of a raw core $N$ of
relative index $d$ that are $H$-isometric to $L$. It is computed by enumerating order-$d$ isotropic subgroups of
$N^\#/N$ for the discriminant bilinear pairing and testing the
isometry type of their inverse images, with the prescribed
norm-integrality convention. Counting the isometries $\sigma\in H$ with $\sigma N\subset L$
by the overparent $\sigma^{-1}L$ gives
$e_H(N,L)|\Gamma_H(L)|$; counting by the image $\sigma N$ gives the
number of raw images of $N$ times $|\Gamma_H(N)|$.
Hence the weight for a raw core is
$|\Gamma_H(N)|/(e_H(N,L)|\Gamma_H(L)|)$.

Integral isometries and automorphism groups can be enumerated
finitely. In $U^t A U=B'$, the $j$th column $v$ satisfies
$v_i^2\le B'_{jj}(A^{-1})_{ii}$ by Cauchy--Schwarz; the remaining
inner products and the determinant complete the isometry test.

\subsection{The local factors}
Let $P_r(p)$ count nonisotropic lines in a nonsingular $r$-dimensional quadratic space over $\mathbb F_p$, with $P_0=0$. For $r>0$ and odd $p$ it equals
\[
 p^{r-1}-
 \begin{cases}
 \left(\frac{(-1)^{r/2}\det B}{p}\right)p^{r/2-1},&r\ \mathrm{even},\\
 0,&r\ \mathrm{odd}.
 \end{cases}
\]
Indeed diagonalize the form and count zeros by additive character orthogonality. Frequency zero gives $p^{r-1}$ affine zeros. At each nonzero frequency $u$ the diagonal Gauss sums have product $\chi(u)^r\chi(\det B)G^r$, where $\chi$ is the quadratic character and $G^2=\chi(-1)p$. The latter identity follows from $\bar G=\chi(-1)G$ and $G\bar G=p$ by additive orthogonality. Sum over $u$, subtract the zero count from $p^r$, then divide by $p-1$. For scalar $h$ the local series is therefore
\begin{equation}\label{eq:goodfactor}
 1+\sum_{e\ge1}g_{p,1}(e)T^e
 =1+\frac{P_n(p)T}{1-p^{n-1}T}.
\end{equation}
There are $p^{(e-1)(n-1)}$ lifts of each line, by fixing one unit coordinate to one. If $p\nmid 2\operatorname{ord}(h)$, a fixed nonisotropic line has eigenvalue $+1$ or $-1$ modulo $p^e$, since its eigenvalue has square one. For an even multiple $k$ of the order prime to $p$, the projectors $k^{-1}\sum_{j=0}^{k-1}(\pm h)^j$ onto these two summands commute with reduction. Their images are nondegenerate and split off orthogonally since the projectors are self-adjoint. The nonconstant part of \eqref{eq:goodfactor} for $h$ is the sum of $P_{r_\pm}(p)T/(1-p^{r_\pm-1}T)$ over its nonzero eigenspaces, using the ranks and determinants of these summands. At every excluded prime the finite congruence test still applies.

\section{A finite incidence formula over copies of the parent}\label{sec:inc}

Suppose $V$ is positive definite and $H$ denotes either
$O(V)(\mathbb Q)$ or $SO(V)(\mathbb Q)$. For $m\ge1$, set
\[
 \mathcal P_m(L)=\{gL:g\in H,\; mgL\subset L\}.
\]
Writing \(g=A/m\), we enumerate integral
matrices \(A\) satisfying \(A^tB_LA=m^2B_L\), using the column
bounds of Section~\ref{sec:core}; for the proper count we retain
positive determinant. For each subset \(F\subset\mathcal P_m(L)\) containing \(L\), put \(K_F=\bigcap_{P\in F}P\), \(d_F=[L:K_F]\). Write \(G_F=\{u\in H:uF=F\}\), computable in the finite group \(\Gamma_H(K_F)\). For \(u\in G_F\) let \(c(u,F)\) be the number of cycles of its permutation of \(F\). Write
\[
 b_{K,u}(k)=\#\{M\subset K:[K:M]=k,\ uM=M\}.
\]

If $u$ permutes a nonempty finite set $\Omega$, the $u$-invariant
subsets are unions of cycles. Hence
\[
 \sum_{\substack{\varnothing\ne F\subseteq\Omega\\uF=F}}
 (-1)^{c(u,F)-1}=1.
\]
\begin{theorem}\label{th:incidence}
At every index and in every positive dimension the actual coefficient is
\[
 a_H(m)=\frac1{|\Gamma_H(L)|}
  \sum_{\substack{L\in F\subset\mathcal P_m(L)\\ d_F\mid m}}
  \frac1{|F|}\sum_{u\in G_F}(-1)^{c(u,F)-1}\, b_{K_F,u}(m/d_F).
\]
Each \(b_{K,u}\) is multiplicative. Its prime-power series is given at every prime by the finite expression in Theorem~\ref{th:matrixlocal}.
\end{theorem}
\begin{proof}
For an actual child \(M\), let \(\mathcal P(M)\) be the parents in \(H L\) containing \(M\). Each has the covolume of \(L\), so \(mP\subset M\) for every such \(P\); the set is finite and nonempty. Apply the invariant-subset identity above to every \(u\in\Gamma_H(M)\), then average over \(\Gamma_H(M)\) and sum over the represented classes \(M\). Equivalently, sum over \(H\)-classes of incident pairs \((M,F)\) with coefficient
\(
 |\Gamma_H(M,F)|^{-1}\sum_{u\in\Gamma_H(M,F)}(-1)^{c(u,F)-1}.
\)

For each invariant summand, mark any vertex \(P\in F\), count all \(|F|\) possible marks, and divide by \(|F|\). Under orbit--stabilizer, summing over marked orbits gives exactly \(|F|\) times the unmarked stabilizer weight. Here \(u\) preserves \(F\), although it may permute its marked vertices. Transporting the marked parent to \(L\) leaves the group
\(\Gamma_H(L)\). We therefore sum over raw pairs \(F,M\),
with weight \(1/|\Gamma_H(L)|\). The conditions become \(M\subseteq K_F\), \([K_F:M]=m/d_F\), and \(uM=M\) with \(uF=F\). Interchanging the finite sums yields the formula. Primary decomposition gives the multiplicativity of \(b_{K,u}\).
\end{proof}

\subsection{The fixed-submodule factor at any prime}
Set \(\Lambda=K_p\), \(R=\mathbb Z_p[u]\), \(E=\mathbb Q_p[u]\). Finite order makes the minimal polynomial squarefree in characteristic zero, so \(E\) is a product of fields \(E_i\). Put \(\mathcal O=\prod_i O_i\), the product of their integer rings. Choose a nonnegative integer \(a\) with \(p^a\mathcal O\subset R\), for example by a lattice index calculation. Write \(\widehat\Lambda=\mathcal O\Lambda\simeq\bigoplus_i O_i^{r_i}\), using a chosen basis on each summand. The first inclusion below follows from \(p^a\mathcal O\subset R\):
\[
 p^a\widehat\Lambda\subset\Lambda\subset\widehat\Lambda .
\]
For \(P\subset\widehat\Lambda\) of full rank write \(v(P)=\log_p[\widehat\Lambda:P]\). Enumerate the finite set
\[
 \mathcal B=\{ P:p^a\widehat\Lambda\subset P\subset\widehat\Lambda,\quad
                        RP=P,\ \mathcal O P=\widehat\Lambda\}
\]
as subgroups of the finite quotient. Fix a uniformizer \(\pi_i\) and a system of residue digits in each \(O_i\). Write \(q_i=p^{f_i}\) for its residue cardinality and \(s_i=v_{\pi_i}(p^a)\). For row \(k=1,\ldots,r_i\) choose \(0\le t_{ik}\le s_i\). Form a block upper triangular matrix \(\bar D\) over \(\mathcal O/p^a\mathcal O\) with diagonal \(\pi_i^{t_{ik}}\) modulo \(p^a\). The \(r_i-k\) entries to its right run independently over the digit representatives of \(O_i/\pi_i^{t_{ik}}O_i\). There are no off-block entries. The containment
\(\bar D P\subset\Lambda\bmod p^a\widehat\Lambda\) is tested in
the finite quotient.

\begin{theorem}\label{th:matrixlocal}
For a finite-order automorphism of any integral lattice, the local series \(\sum_{e\ge0}b_{K,u}(p^e)T^e\) equals
\[
 \sum_{(t_{ik})}
 \left(\prod_{i,k}T^{f_i t_{ik}}\right)
 \left(\prod_{t_{ik}=s_i}(1-q_i^{r_i-k}T^{f_i})^{-1}\right)
 \sum_{\bar D}\ 
 \sum_{\substack{P\in\mathcal B\\\bar D P\subset\Lambda\bmod p^a\widehat\Lambda}}
 T^{v(P)-v(\Lambda)}.
\]
In particular its denominator divides
\(\prod_i\prod_{j=0}^{r_i-1}(1-q_i^jT^{f_i})\) after cancellation. If \(R=\mathcal O\), its numerator can be taken to be one. This applies at \(p\nmid\operatorname{ord}(u)\), since \(X^{\operatorname{ord}(u)}-1\) is separable modulo \(p\).
\end{theorem}
\begin{proof}
For an \(R\)-stable \(M\subset\Lambda\) let \(\widehat M=\mathcal O M\). Each component of \(\widehat M\) has a unique upper triangular Hermite matrix with diagonal \(\pi_i^{l_{ik}}\) and entries to the right in the fixed digit representatives modulo \(\pi_i^{l_{ik}}\). This follows by projecting successively on the last coordinate and reducing each lifted basis vector by the columns already chosen. The combined matrix \(D\) identifies \(\widehat\Lambda\) with \(\widehat M\). Now \(p^a\widehat M\subset M\), so \(P=D^{-1}M\) belongs to \(\mathcal B\). Conversely every pair \(D,P\) with \(DP\subset\Lambda\) gives such an \(M\), and \(\mathcal O(DP)=D\widehat\Lambda\) recovers \(D\). The containment test depends only on \(D\bmod p^a\), while the
index exponent is \(\sum f_i l_{ik}+v(P)-v(\Lambda)\).
Clip \(l_{ik}\) to \(t_{ik}=\min(l_{ik},s_i)\). For \(t_{ik}<s_i\) the residue matrix fixes all the entries of that row. For \(l_{ik}\ge s_i\) every admissible residue row has exactly \(q_i^{(l_{ik}-s_i)(r_i-k)}\) lifts. Summing these multiplicities with their index weights gives the displayed geometric factors. Any apparent negative power from \(v(P)-v(\Lambda)\) is offset by the index of the residue matrix lifted with \(l_{ik}=t_{ik}\), because its image of \(P\) lies in \(\Lambda\). If \(a=0\) the only lattice \(P\) is \(\widehat\Lambda=\Lambda\) and all \(t_{ik}=0\). At a prime away from the order of \(u\), factor the separable reduction and lift the distinct factors by Hensel's lemma. The resulting finite unramified integer rings are integrally closed and the summands corresponding to the minimal polynomial give \(R\).
\end{proof}
For invariant-submodule zeta functions of orders, see Solomon
\cite{Solomon} and Rossmann's local formulas and exceptional-prime
corrections \cite[Thm.~A and Rem.~5.6]{Rossmann}. Together with the incidence formula, these local fixed-submodule
series determine the intrinsic coefficients.


\section{Meromorphic continuation and the counting law}\label{sec:analytic}
Set $S_d=S\cup\{p:p\mid d\}$, and retain the nonisotropic-line
count $P_n$ of Section~\ref{sec:core}. Define
\[
 A_d(s)=\prod_{p\in S_d}(1-p^{n-1-s})
         \prod_{p\notin S_d}\bigl(1+(P_n(p)-p^{n-1})p^{-s}\bigr),\qquad
 W_H(d)=\sum_{[N]\in\mathcal C_H(d)}\frac{c_H}{|\Gamma_H(N)|}.
\]
\begin{theorem}\label{th:analytic}
For every positive definite $L$ of rank $n\ge2$, each orientation convention has
\[
 Z_H(s)=\zeta(s-n+1)D_H(s)+E_H(s),\qquad
 D_H(s)=\sum_{d\ge1}W_H(d)d^{-s}A_d(s),
\]
where $D_H,E_H$ are absolutely convergent Dirichlet series on $\Re s>n-1$. The abscissa of convergence is $n$, and the unique pole on $\Re s>n-1$ is simple with residue $R_H=D_H(n)>0$. For real $X\ge2$ and every $\epsilon>0$,
\[
 \sum_{m\le X}a_H(m)=\frac{R_H}{n}X^n+O_{L,\epsilon}(X^{n-1+\epsilon}).
\]
\end{theorem}
\begin{proof}
\emph{The core majorant.}
Let $B(d)$ be the number of raw cores. Primary decomposition of subgroups of the finite module $L/dL$ gives multiplicativity of $B$. At $p\in S$, its local factors have coefficients $b_n(p^e)$. Elsewhere they have coefficients $b_n(p^e)-P_n(p)p^{(e-1)(n-1)}$ for $e\ge1$. Uniformly in $e\ge1$,
\[
 p^{(n-1)e}\le b_n(p^e)\le p^{(n-1)e}\prod_{j=1}^{n-1}(1-p^{-j})^{-1}
 \le p^{(n-1)e}(1+C_n/p).
\]
The first two inequalities follow by extracting the highest-weight
geometric factor from the HNF series. Since $P_n(p)=p^{n-1}+O_n(p^{n-2})$, the bad coefficients outside $S$ are $O_n(p^{(n-1)e-1})$. For $\sigma=n-1+\delta$, $\delta>0$, their weighted local sum is $O_{n,\delta}(p^{-1-\delta})$. Hence $\sum B(d)d^{-\sigma}<\infty$.

\emph{Scalar contributions.}
For scalar $h$, \eqref{eq:goodfactor} supplies exactly $\zeta(s-n+1)A_d(s)$. The sum of absolute values of coefficients of $A_d$ weighted at $\sigma$ is at most
\[
 C_{n,\delta}\prod_{p\mid d}(1+p^{-\delta})\ll_{n,\delta,\eta,S} d^\eta\qquad(\eta>0).
\]
Absolute convergence of the product outside $S_d$ follows from
$\sum_p p^{-1-\delta}<\infty$; moreover,
$1+p^{-\delta}\le p^\eta$ for all sufficiently large $p$.
With $W_H(d)\le B(d)$ and $0<\eta<\delta$, the Dirichlet series for $D_H$ converges absolutely.

\emph{Nonscalar contributions.}
For a nonscalar $h\in\Gamma_H(N)$, reduction modulo any odd prime cannot be scalar. Otherwise conjugation by $h$ on $\mathrm{End}_{\mathbb Z}(N)$ would be trivial modulo $p$. A finite-order integral matrix congruent to the identity modulo an odd prime is the identity: a nontrivial prime-order power has the form $1+p^a C$, $a\ge1$, with some entry of $C$ a unit. By the binomial expansion its $\ell$th power minus one has minimum entry valuation $a$ for $\ell\ne p$ and $a+1$ for $\ell=p$. Trivial conjugation forces $h$ to be scalar over $\mathbb Q$.

Positive definiteness makes $\Gamma_H(N)$ finite. A nonscalar matrix over $\mathbb F_p$ has at most $n$ eigenspaces for rational eigenvalues, each of dimension at most $n-1$. It therefore fixes $O_n(p^{n-2})$ projective lines. Bounding their lifts gives, uniformly,
\[
 g_{p,h}(e)\ll_n p^{(n-1)e-1}\quad(e\ge1,\ p\notin S_d).
\]
At $\sigma$ these local products are bounded in terms of
$n,\delta$, independently of nonscalar $h$. Averaging over
$\Gamma_H(N)$ and using the core majorant $B(d)$ proves absolute
convergence of $E_H$. The HNF majorant justifies the rearrangements on $\Re s>n$, and the convergent expressions extend the identity to $\Re s>n-1$.

\emph{The residue and the summatory law.}
Each $A_d(n)$ is positive: each of its factors outside $S_d$ is at least $1-p^{-1}$ and the product converges absolutely, and the excluded factors equal $1-p^{-1}$. Also $W_H(1)>0$. For the summatory formula let $u(k)$ be the coefficients of $D_H$. In the convolution with $j^{n-1}$ use $\sum_{j\le y}j^{n-1}=y^n/n+O_n(y^{n-1})$ for $y\ge1$. The main-term tail is bounded by $X^{n-1+\eta}\sum_k|u(k)|k^{-(n-1+\eta)}$ for $0<\eta<\min(1,\epsilon)$. Absolute convergence gives the same bound for the power-sum
error and the summatory function of $E_H$. Since the main coefficient
is positive, partial summation gives divergence at $n$.
\end{proof}

\begin{corollary}\label{cor:KLL}
For each scalar $\rho\in H$ the transporter-weighted term of Kim--Lee--Lee has a simple pole at $n$ of residue $R_H/c_H$ and is meromorphic on $\Re s>n-1$. The sum over \emph{all nonscalars} converges absolutely there. For even $n$ the entire determinant-minus-one sum is holomorphic there. For odd $n$ that sum equals $\tfrac12 Z_O$.
\end{corollary}
\begin{proof}
In the notation of \cite[\S2.4, (2.4)--(2.5)]{KLLpre}, a raw $M$ fixed by $\rho$ contributes the reciprocal of the number of $\sigma\in H$ with $\sigma M\subset L$. At a fixed index this denominator is the number of images times $|\Gamma_H(M)|$. Thus the contribution of a scalar equals $\sum_{[M]}1/|\Gamma_H(M)|$. In each restored-core orbit of children the stabilizer is exactly $\Gamma_H(M)$, by uniqueness of restoration. The identity term of Burnside evaluates that same sum, so the total scalar part is precisely $\zeta(s-n+1)D_H(s)$. The remaining transporter contributions have nonnegative
coefficients and sum to $E_H$. In odd rank $\rho\mapsto-\rho$ pairs determinant signs with identical fixed sets and denominators.
\end{proof}

\section{Orientation, prime extensions and Euler products}\label{sec:orientation}
For an ordinary represented class let its two possible orientations be identified if it has an improper automorphism. If such an automorphism exists, the class is called achiral. A chiral class can have either one or two orientations represented in $L$. The set of represented proper classes need not be invariant under orientation reversal. If $S_m$ is that set, with involution $\iota$ acting on the full ambient set of proper classes, then
\[
 2a_O(m)=|S_m\cup\iota S_m|+|S_m\cap\operatorname{Fix}(\iota)|.
\]

\begin{theorem}\label{th:orientation}
In every even rank $n\ge2$,
\[
 \tfrac12 R_{SO}\le R_O<R_{SO},\qquad
 R_O=\tfrac12 R_{SO}\ \Longleftrightarrow\ L\text{ has an improper automorphism}.
\]
In particular the excess of proper over ordinary classes has a positive leading term in its summatory function, even when $L$ itself has no improper automorphism.
\end{theorem}
\begin{proof}
Group the proper core masses at each $d$ by ordinary core class. Both scalar groups have size two. An achiral ordinary core contributes twice its ordinary mass to the proper mass, since the stabilizer loses a factor two. A chiral core contributes its ordinary mass times the number of orientations actually represented. The factors $A_d(n)$ depend only on $d,L,S$. If $L$ has an improper automorphism, it supplies the other orientation of every chiral core. If $L$ is chiral the index-one class occurs in just one orientation, proving strict inequality on that side. For the other strict inequality choose a rational orthogonal basis and clear denominators to obtain an achiral diagonal sublattice $K\subset L$ of index $b$. Enlarge $S$ by the primes dividing $b$. The decomposition is valid for this finite set and the residues of the same functions are unchanged. This $K$ is now a core and contributes strictly more proper than ordinary mass. Apply Theorem~\ref{th:analytic}.
\end{proof}

Put $u_H(M)=|\Gamma_H(M)|/c_H$ and $\mathcal M_H(b)=\{[M]_H:[L:M]=b,\ M\subset L\}$.
\begin{lemma}\label{lem:pr}
For fixed $b\ge1$ and primes $p\to\infty$,
\[
 a_H(bp)=p^{n-1}\sum_{[M]\in\mathcal M_H(b)}u_H(M)^{-1}+O_{L,b}(p^{n-2}).
\]
\end{lemma}
\begin{proof}
Include all divisors of $b$ that are primes in $S$. When $p\nmid b$ is regular, saturation gives the actual classes at index $b$. Equations \eqref{eq:goodfactor} and the nonscalar fixed-line estimate give the stated leading term. At an irregular $p$ the number of raw choices is at most $b_n(b)(b_n(p)-P_n(p))=O_{b,n}(p^{n-2})$ by primary decomposition. The same bound holds for the corresponding classes.
\end{proof}

\begin{theorem}\label{th:nonmult}
For every positive definite rational lattice of rank $n\ge3$, neither orientation count is multiplicative. Neither Dirichlet series admits the normalized Euler product of its prime-power subseries.
\end{theorem}
\begin{proof}
Write $u=u_H(L)$. Multiplicativity and Lemma~\ref{lem:pr} at $b$ and $1$ would give
\[
 \sum_{[M]\in\mathcal M_H(b)}u_H(M)^{-1}=a_H(b)/u\qquad\text{for every }b.
\]
For $b=\ell$ a large regular-eligible prime, this weighted count is $\ell^{n-1}/u+O_L(\ell^{n-2})$. Indeed its regular contribution is exactly $P_n(\ell)/u$ by orbit--stabilizer and the irregular part is bounded by its raw count. The unweighted $a_H(\ell)$ also equals $\ell^{n-1}/u+O_L(\ell^{n-2})$ by the lemma. Letting $\ell\to\infty$ forces $u=1$. The displayed equality now forces $u_H(M)=1$ for every represented sublattice. A rational orthogonal basis supplies a diagonal sublattice inside $L$. Changing two of its coordinate signs is a proper nonscalar isometry in every $n\ge3$, a contradiction for either $H$. Uniqueness of Dirichlet coefficients gives the Euler-product assertion.
\end{proof}

\section{The scalar genus mass and its Euler product}\label{sec:genusmass}
Keep the positive definite hypotheses and fix $H$. Consider the finite set of all $H$-classes of parents in the genus of $L$ in $V$, with representatives $L_1,\ldots,L_h$ and mass $\mathfrak m_H=\sum_i|\Gamma_H(L_i)|^{-1}$. Local ordinary and proper types coincide by the integral reflection
argument of Section~\ref{sec:indef}; the fixed orientation is used at
the real place. The genus is taken in the fixed rational space $V$.

For a local sublattice $P\subset L_p$ of index $p^e$, let $b(P,L_p)$ count its overlattices isometric to $L_p$ of relative index $p^e$. This finite positive number is obtained by testing isotropic subgroups of the local discriminant module. Put
\[
 t_p(e)=\sum_{\substack{P\subset L_p\\{}[L_p:P]=p^e}}b(P,L_p)^{-1}.
\]
Let $\mathcal U_H(m)$ be the union of the sets of represented child classes at index $m$ over all the $L_i$. Then
\begin{equation}\label{eq:genusmass}
 \sum_{[M]\in\mathcal U_H(m)}\frac{1}{|\Gamma_H(M)|}
       =\mathfrak m_H\prod_{p^e\Vert m}t_p(e).
\end{equation}
In particular for a parent of $H$-class number one in its genus its entire scalar contribution has an Euler product. Its residue is explicitly
\[
 R_H=\frac{c_H}{|\Gamma_H(L)|}
        \prod_p\left[(1-p^{-1})\sum_{e\ge0}t_p(e)p^{-ne}\right].
\]
A parent above $M$ in this genus corresponds to a collection of local overlattices of the prescribed types. Patching is bijective; every result is one of the global genus classes. Hence their total number is $b(M)=\prod_p b(M_p,L_p)$. Counting transporters as in Theorem~\ref{th:core}, the number of raw copies of $M$ in $L_i$ divided by $|\Gamma_H(L_i)|$ equals the number of its overparents of class $i$ divided by $|\Gamma_H(M)|$. Sum over $i$, divide by $b(M)$, then interchange the sum over raw sublattices and the finite parent sum. The local raw sum of $1/b(M)$ is identical for every $L_i$ by local isometry, and factors by primary decomposition. This proves \eqref{eq:genusmass}.

At $p\notin S$, regular kernels have $b(P,L_p)=1$, and
irregular kernels contribute weight at most one. The estimates in Theorem~\ref{th:analytic} imply that after extracting $\zeta(s-n+1)$ the normalized product is absolutely convergent on $\Re s>n-1$. At $s=n$ each factor is strictly positive. The other transporter terms are holomorphic there, giving the claimed residue for a single-class genus. 

\section{The growth of classes admitting an improper symmetry}
In even positive rank $n\ge4$ let $b^-(m)$ count the ordinary represented classes whose intrinsic automorphism group contains an element of determinant $-1$; no requirement is made that it fix the parent. Then for $X\ge2$, as $X$ tends to infinity,
\begin{equation}\label{eq:reflection-growth}
 X^{n-1}\ll_L\sum_{m\le X}b^-(m)\ll_{L,\epsilon}X^{n-1+\epsilon}.
\end{equation}
The series with coefficients $b^-(m)$ has abscissa $n-1$ and diverges there. Thus the absolute-convergence half-plane for the improper term of Corollary~\ref{cor:KLL} is sharp.

The coefficient of the improper transporter sum is $b^-(m)/2$,
since an intrinsic automorphism group containing an improper element
has equally many automorphisms of each determinant sign. The upper
bound follows. For a lower bound take a diagonal orthogonal sublattice $K$ of index $b$ in $L$, and place all divisors of $2b\det B_K$ that are prime in $S$. Among its regular children take those fixed by the sign change in its last coordinate. At a prime $p\notin S$ take all primitive projective normals supported on the first $r=n-1$ coordinates and nonisotropic modulo $p$. They give distinct raw kernels in $K_p$ by the unit-coordinate chart. Since $r$ is odd, their number at $p^e$ is exactly $p^{(r-1)e}$, by the finite-field calculation above. For every $t$ prime to $S$, there are therefore at least $t^{n-2}$ such children of index $bt$ in $L$. A rational isometry between two of them stabilizes their restored lattice $K$. Dividing by $|\Gamma_O(K)|$ bounds the number of represented ordinary classes from below. Summing over $t\le X/b$ prime to $S$ gives the lower bound by finite inclusion--exclusion, and the subseries at $n-1$ dominates a positive multiple of the harmonic series on these $t$.

\subsection*{The distribution of represented orientations}
Suppose \(n\ge2\) is even and \(L\) is positive definite. Let \(U(m)\) and \(T(m)\) count chiral ordinary classes represented with exactly one and exactly two proper orientations, respectively. Then
\[
\begin{aligned}
 \sum_{m\le X}U(m)&=\frac{2R_O-R_{SO}}{n}X^n+
                         O_{L,\epsilon}(X^{n-1+\epsilon}),\\
 \sum_{m\le X}T(m)&=\frac{R_{SO}-R_O}{n}X^n+
                         O_{L,\epsilon}(X^{n-1+\epsilon}).
\end{aligned}
\]
The two coefficient identities are \(a_O=b^-+U+T\) and
\(a_{SO}=b^-+U+2T\). In every even dimension the improper transporter coefficient is \(b^-/2\), including in rank two, so the absolute convergence of the nonscalar remainder in Corollary~\ref{cor:KLL} bounds its summatory function by \(O_{L,\epsilon}(X^{n-1+\epsilon})\). Apply Theorem~\ref{th:analytic}. In particular a chiral parent represents a positive limiting proportion of each kind among all ordinary classes ordered by index. For an achiral parent \(U(m)=0\) for every \(m\).

For either orientation group in any positive rank \(n\ge2\), all but \(O_{L,\epsilon}(X^{n-1+\epsilon})\) represented classes of index at most \(X\) have \(\Gamma_H(M)=H\cap\{1,-1\}\). Indeed each exceptional class contributes \(1-c_H/|\Gamma_H(M)|\ge1/2\) to the nonscalar transporter sum. For ordinary equivalence such classes therefore have zero limiting proportion.


\section{All indices for the cubic ternary lattice}\label{sec:cubic}
Take $L=\mathbb Z^3$ with the standard dot product. Write
$a(m)$ for the number of intrinsic classes at index $m$ and
$r_3(b)=\#\{(x,y,z)\in\mathbb Z^3:x^2+y^2+z^2=b\}$.
Ordinary and proper classes coincide in rank three. The coefficients
are expressed in terms of $r_3$, quadratic residue symbols and rational
functions of one variable.
For an odd prime $p$ write $\chi_p(b)=(b/p)$. In the following formulas
$T$ is formal; products over $p^k\parallel m$ use only $k>0$.

\paragraph{Local factors at odd primes.}
Put
\[
 F_1=\frac1{(1-T)(1-pT)(1-p^2T)},\quad
 F_2=\frac1{(1-T)^2(1-pT)}.
\]
Also define
\[
 J_p(b)=\frac1{1-T}
 \begin{cases}(1-T)^{-2}&\chi_p(b)=1,\\
 (1-T^2)^{-1}&\chi_p(b)=-1.\end{cases}
\]
The identity row is $F_1(1-(p+1)T/2)$. The half-turn row has an
axis parameter $b$ and a squarefree odd reversal index $e$. At $p\mid e$
it is $TF_2/2$. At $p\nmid e$ it is
$F_2(1-(1+\chi_p(-b))T/2)$. Here $b=e$ or $2e$.
The quarter-turn row is $J_p(-1)(1-(1+\chi_p(-1))T/2)$.
The third-turn and sixth-turn rows at $p\ne3$ are both
$J_p(-3)(1-(1+\chi_p(-3))T/2)$. At $p=3$ they are, respectively,
\[
 1+\frac{T(1+2T)}{2(1-T)^2},\qquad \frac{T}{2(1-T)^2}.
\]
The sixth-turn row is used only when $3\mid m$.

\paragraph{Local factors at two.}
For the identity retain $F_1$ with $p=2$ and set
$P=(1+3T)/((1-T)^2(1-2T))$. The other rows are
\[
\begin{array}{c|c}
 \text{rotation}&\text{factor at }2\\ \hline
 \text{half-turn, }b\text{ odd}&P\\
 \text{half-turn, }b\text{ even}&(1-2T)^{-1}+TP\\
 \text{quarter-turn}&(1+T)/(1-T)^2\\
 \text{third- or sixth-turn}&((1-T)(1-T^2))^{-1}.
\end{array}
\]
For each row $j$ write $W_j(m)=\prod_{p^k\parallel m}[T^k]A_{j,p}(T)$.
Use $W_2(b,e;m)$ for the half-turn row.
\begin{theorem}\label{thm:cubic-all}
For every $m\ge1$,
\[
\begin{split}
24a(m)={}&W_1(m)+6W_4(m)+8W_3(m)
  +8\mathbf1_{3\mid m}W_6(m)\\
 &+\sum_{\substack{e\mid m\\e\ {\rm odd\ squarefree}}}
  \left(\frac{r_3(e)}2W_2(e,e;m)
       +\frac{r_3(2e)}2W_2(2e,e;m)\right).
\end{split}
\]
Here the subscripts $1,4,3,6$ denote rotation orders.
\end{theorem}
In particular $a(2)=a(3)=3$, $a(6)=13$. The cubic Dirichlet series
does not have multiplicative coefficients. For $p\ge5$ the rows give
\[
 a(p)=\frac{2p^2+19p+70+3\chi_p(-1)-6\chi_p(-2)+8\chi_p(-3)}{48}
       +\frac{r_3(p)+r_3(2p)}{96}.
\]
At any fixed prime $p$, the formula gives a rational series in $T$
for the coefficients $a(p^k)$: only $e=1,p$ occur, and at $p=2$
only $e=1$.

\subsection*{Maximal parents and their local trees}
An integral unimodular positive rank-three lattice is isometric to $L$.
The open ball of radius $\sqrt2$ has volume greater than $2^3$,
so Minkowski's theorem gives a norm-one vector. Its orthogonal complement
is integral unimodular; the open disk of radius $\sqrt2$ has area greater
than $2^2$, giving another norm-one vector. Thus every integral self-dual
lattice in the rational cubic space is a possible parent. The positive
integral stabilizer $\Gamma=\mathrm{SO}(L)$ consists of the 24 determinant-one
signed permutations.

Over $\mathbb Q_2$ the norm-integral vectors are exactly
$L_2=\mathbb Z_2^3$: an integral triple with an odd coordinate has sum
of squares nonzero modulo four. Consequently the self-dual parent at two
is unique. For $p$ odd let the vertices be the self-dual lattices in
the given local space; join vertices whenever their intersection has
index $p$ in both. For any two vertices $P,R$ the elementary divisor
exponents relative to one another are $(-h,0,h)$ since duality negates
these exponents. An index-$p$ hyperplane $K\subset P$ corresponds to
a line $\ell=(K/pP)^\perp$ in the nonsingular space $P/pP$. If
$\ell$ is anisotropic the discriminant group $K^*/K$ is cyclic of
order $p^2$: split off a unit vector lifting $\ell$, then $K$ replaces
this summand by its $p$-multiple. There is just one index-$p$ overgroup.
If $\ell$ is isotropic, choose a lift $v\in P$ with $v^2\in p^2\mathbb Z_p$
by adjusting $v$ by $p$ times a vector pairing to a unit with it.
Then $K+\mathbb Z_p(v/p)$ is a second self-dual parent. The two
distinct order-$p$ subgroups of $K^*/K$ defining these parents show
that the group is elementary; its pairing, multiplied by $p$, is a
nondegenerate hyperbolic plane. Its two isotropic lines are all the
possibilities. The neighbors therefore correspond to the isotropic lines. There are $p+1$ lines, since a nondegenerate ternary
form over a finite field of odd order has an isotropic point and the
conic through that point is parametrized by the projective line.
After diagonalization and fixing the third coordinate, isotropy follows
by intersecting two subsets of $\mathbb F_p$, each containing $(p+1)/2$
quadratic values.

Choose a Smith basis of $P$ in which $R$ has generators
$p^h u,v,p^{-h}w$ with $h>0$. The lattice with generators
$pu,v,p^{-1}w$ is integral by the pairings in $P$ and $R$ and has
unit determinant. Continuing with each exponent between 0 and $h$
gives a path of length $h$; its members all contain $P\cap R$.
The exponent $h$ is the least integer such that $p^hR\subset P$,
so it changes by at most one along edges. The line $p^hR\bmod pP$
is spanned by $w$ and is isotropic. Any other isotropic line has
nonzero pairing with $w$, since a nonsingular ternary space contains
no totally isotropic plane. The corresponding neighbor omits $p^hR$,
so its distance exponent to $R$ is $h+1$. Every vertex has a unique edge toward the root, with no edges
within a level; hence the neighbor graph is a tree. For a fixed full lattice $M_p$, the set of parents containing $M_p$
is a finite convex subtree when nonempty. Every vertex on the path
between two such parents contains their intersection, and every parent
lies between $M_p$ and $M_p^*$.

For a represented lattice $M$ of determinant $m^2$ the product of these
subtrees is the complex of its possible parents. Each vertex gives an
actual global parent by intersecting its local lattices with the rational
space. This local-to-global correspondence follows by placing
all lattices between $cL$ and $c^{-1}L$ and decomposing their finite
quotient into its primary groups. The dyadic factor and every $p\nmid m$
contribute one vertex. The index of $M$ in every parent is $m$.

\subsection*{The signed count of marked cubes}
The product of parent trees is contractible. For $g\in\mathrm{SO}(M)$
the alternating sum of its traces on the rational cellular chains is 1:
cancel the traces on boundaries in successive dimensions to leave
only $H_0$. On a preserved $r$-cube the contribution is
$(-1)^r(-1)^{b_g}$ where $b_g$ is the number of reversed local edges.
Prime places cannot be permuted. Averaging over $\mathrm{SO}(M)$
and summing over represented classes, we may group terms by unanchored
cubes. For such a cube $C$, let $K_C$ be the intersection of its
vertices and $S_C$ its rational proper stabilizer. It is finite, since the stabilizer of any one
of its finitely many vertices is finite. If edges occur at the primes
of $d$, its intersection has index $d$ in every vertex. The lattices
with $C$ among their parents are exactly the index-$m/d$ sublattices
of $K_C$. Changing the order of the finite incidence sum, the contribution
of the orbit of $C$ is
\[
 \frac{(-1)^{\omega(d)}}{|S_C|}
 \sum_{g\in S_C}(-1)^{b_g}\,
 \#\{M\subset K_C:[K_C:M]=m/d,\ gM=M\}.
\]
In both orbit decompositions, a marked pair has stabilizer
$S_C\cap\mathrm{SO}(M)$, on which the character $(-1)^{b_g}$ is
evaluated. We may mark one vertex of each cube as $L$.
Summing over all cubes anchored there replaces the preceding denominator
by $2^{\omega(d)}|\Gamma|$. To see this, the sum over $S_C$-orbits
of vertices of $1/|(S_C)_P|$ equals $2^{\omega(d)}/|S_C|$;
after identifying $P$ with $L$ its remaining changes of coordinates
are precisely $\Gamma$, independently of whether the cube admits an
edge-reversing symmetry.

For an anchored cube $gL$ is a corner and $g^2L=L$ because $g$ preserves
or reverses each local edge. Write $e$ for the product of primes where
$gL_p\ne L_p$. At these primes an edge is forced, $K_{C,p}=L_p\cap gL_p$,
and the signs from dimension and orientation cancel. At an odd prime
outside $e$, either there is no edge or one chooses a $g$-fixed isotropic
line at $L_p$. Thus, writing $F_{g,P}(T)$ for the series of $g$-stable
sublattices of $P$ by index exponent, the contribution of $g$ factors
into local coefficients. At a reversed prime the factor is
$TF_{g,L_p\cap gL_p}/2$. At another odd prime it is
\[
 F_{g,L_p}-\frac T2\sum_{\substack{\ell\ \text{isotropic}\\g\ell=\ell}} F_{g,K_\ell}.
\]
At two it is $F_{g,L_2}$. Conversely any rational rotation with square
fixing $L$, with squarefree odd displacement $e\mid m$, and with the
indicated choices at the remaining primes stabilizes exactly that cube.
The product for invariant sublattices follows by primary decomposition
of $K_C/(m/d)K_C$.

\subsection*{Classification of rational rotations}
Write a rational rotation as conjugation on the trace-zero Hamilton
quaternions by a primitive integral quaternion $q=(w,x,y,z)$, unique
up to sign. To see surjectivity, an integral or rational orthogonal
rotation of trace $-1$ is a half-turn about a rational axis, represented
by its pure axis vector. Otherwise its Cayley transform
$(g-I)(g+I)^{-1}$ is rational skew-symmetric and is the cross-product
matrix of a rational vector $v$. The quaternion $(1,v)$ gives back
$g$ by multiplying out the Cayley transform. Central scalars are exactly
the kernel. The least denominator of the rotation matrix is the odd part
of $q\bar q$. Indeed its entries have this norm as common denominator,
with diagonal numerators $w^2+x^2-y^2-z^2$ and their cyclic translates
and off-diagonal numerators $2(xy\pm wz)$ and their cyclic translates.
At an odd prime dividing the norm these cannot all vanish: vanishing
forces the six pairwise products of the four coordinates to vanish
modulo $p$, and the zero norm then forces every coordinate to vanish.
At two the primitive norm has valuation at most two. If it has valuation
one all numerators are even; if it has valuation two all four coordinates
are odd, and all numerators are divisible by four. Because of the
symmetric three Smith exponents the least denominator is also
$[L:L\cap gL]$. This agrees with the cubic coincidence index \cite{Baake}.

The denominator here is squarefree and equals $e$.
As $q^2=2wq-q\bar q$, its coefficients are divisible by $p\mid e$
if and only if $p\mid w$. The square rotation is integral precisely
when this holds for every such $p$: divide $q^2$ by $p$ to remove the
entire $p$-part of its norm. For $e\ge5$ the estimate
$|w|\le2\sqrt e<e$ forces $w=0$. Pure primitive quaternions have
norm $e$ or $2e$, since a sum of three squares divisible by four
has all coordinates even. Conversely every triple of norm $e$ or
$2e$ is primitive by squarefreeness. These give $r_3(e)/2,r_3(2e)/2$
rotations. For $e=3$ there are also the eight rotations represented
by $(3,\pm1,\pm1,\pm1)$, of order six. For $e=1$ there are the 24
rotations of the cube: the identity, nine half-turns (already included
by the pure quaternions), six quarter-turns, and eight third-turns.

\subsection*{Evaluation of the local traces}
At an odd prime the identity fixes all $p+1$ conic points. A half-turn
with primitive axis of norm $b$ at $p\nmid 2b$ has an anisotropic
fixed axis; the determinant of the orthogonal plane has square class
$b$. Hence $1+\chi_p(-b)$ isotropic lines are fixed. For a quarter-turn,
or a third- or sixth-turn at $p>3$, the axis is again anisotropic.
The other two eigenlines are isotropic whenever defined since their
eigenvalues square to a number different from one. The counts are
$1+\chi_p(-1)$ and $1+\chi_p(-3)$ respectively. The quarter-turn calculation also applies at $p=3$. An integral third-turn is conjugate by a signed permutation
to the cyclic permutation of the three coordinates. It fixes the unique eigenline $(1,1,1)$ modulo three,
which is isotropic.

At primes not dividing the rotation order the action decomposes by
the distinct irreducible residue factors of the minimal polynomial:
the Bezout projections lift over $\mathbb Z_p$. For a half-turn the
summands have ranks one and two, and the unrestricted subgroup series
on these summands give $F_2$. For the other rotations with two non-real
eigenvalues, the nontrivial plane either splits into two lines or is
free of rank one over the unramified quadratic integer ring. Its stable
sublattices then give $(1-T)^{-2}$ or $(1-T^2)^{-1}$. This description
applies equally to each stable edge intersection. For the identity use the
rank-three HNF series $F_1$. For the exceptional order-six rotations at
three the polynomial $X-1$ is coprime to $X^2-X+1$ over $\mathbb Z_3$.
The plane is rank one over $\mathbb Z_3[\zeta_6]$, a discrete valuation
ring with residue field of size three: $\zeta_6+1$ is a uniformizer
by its Eisenstein polynomial at three. Thus its factor, including the
axis, is $(1-T)^{-2}$, also on the reversed edge.

At the exceptional primes, the local module calculation reduces
to a fiber-product order. Let $O$ be a local integer ring with
residue field $\mathbb F_p$ and uniformizer $\pi$. Consider the order
\[
 A=\{(x,y)\in\mathbb Z_p\times O:x\bmod p=y\bmod\pi\}.
\]
For an $A$-stable lattice $E$ with saturated module
$(\mathbb Z_p\times O)E=D=p^a\mathbb Z_p\oplus\pi^bO$, one has
$(p,\pi)D\subset E$ (here $(p,\pi)D$ denotes the direct sum of the
two maximal-ideal multiples). The image of $E$ modulo this direct sum
projects onto both residue spaces. The possibilities are the whole
plane or one of the $p-1$ graphs of nonzero maps. Conversely all these
possibilities work. For ambient $\mathbb Z_p\oplus O$ the series is
$F_{\rm split}=(1+(p-1)T)/(1-T)^2$. For ambient $A$, when $a,b\ge1$,
all graphs and the plane occur, with one power of $T$ removed because
the ambient has index $p$. When only one of $a,b$ is zero, the
residue-matching condition is impossible; at $(0,0)$ only the graph
defining $A$ occurs. Thus $F_{\rm glued}=1+TF_{\rm split}$.
For the third-turn at three, evaluation at $(1,\zeta_3)$ identifies
the permutation lattice $\mathbb Z_3[X]/(X^3-1)$ with $A$ for
$O=\mathbb Z_3[\zeta_3]$ and $\pi=1-\zeta_3$.
The two factors have resultant $3$ and common residue field
$\mathbb F_3$; polynomial representatives therefore give pairs
with the same residue. The hyperplane for the fixed isotropic line imposes
coordinate sum zero modulo three, so its module is $3\mathbb Z_3
\oplus\pi O$, isomorphic to the split module. This gives
$F_{\rm glued}-TF_{\rm split}/2$. At two the quarter-turn is
the direct sum of its fixed axis and $\mathbb Z_2[i]$; evaluation
at $(1,i)$ gives the same order $A$ with $\pi=1-i$ owing to the
resultant $2$, and the ambient is split. The remaining third- and
sixth-turns at two have a separate axis and the unramified quadratic
plane.

Finally, a dyadic half-turn has saturation
$R=\mathbb Z_2\oplus\mathbb Z_2^2$ under the two rational eigenprojections.
In coordinates on its projected summands the original lattice is $R$
or the graph of a nonzero linear form on the negative residue plane:
it contains $2R$, the two projections are surjective, and the only
other subspace with surjective projections is the full residue space.
The split case is equivalent to integral eigenprojections, or
$g\equiv I\pmod2$. For a primitive axis $v$ with norm $b$ the half-turn
is $2vv^t/b-I$, so this is exactly the case of odd $b$.
A sublattice with saturation $2^a\mathbb Z_2\oplus D$ for
$D\subset\mathbb Z_2^2$ has four choices (one full residue space and
three graphs), giving the series $P$ for a split ambient.
In a glued ambient, when $a\ge1$ one must have
$D\subset\ker(\ell\bmod2)$ for its defining form $\ell$, and then
every choice works, contributing $TP$. When $a=0$, $D$ must avoid
containment in this kernel; exactly one graph works. The series of
such $D$ is $(1-T)/((1-T)(1-2T))$. The asserted dyadic factors follow.

\subsection*{A branch point inside the domain of continuation}
Let $Z_{\rm cub}(s)=\sum_m a(m)m^{-s}$ and let all powers on the real
line initially have their positive values. Define
\[
 R(s)=(1-2^{1-s})^{-1/2}
 \prod_{p>2}\big(1-\tfrac12(p+1)p^{-s}\big)
                  (1-p^{1-s})^{-1/2}.
\]
For $\Re s>3/2$ the product is nonvanishing and holomorphic:
on compact subsets the logarithm of the $p$-factor is
$O(p^{-\Re s}+p^{2-2\Re s})$, and each factor is nonzero.
For a positive integer $b$ put $P_b(X)=1+3X$ if $b$ is odd and
$P_b(X)=1-X+4X^2$ if $b$ is even. Set
\[
 B(s)=\sum_{e\ {\rm odd\ squarefree}}
 \frac{e^{-s}}{2^{\omega(e)}}\sum_{b\in\{e,2e\}}r_3(b)P_b(2^{-s})
  \prod_{p\nmid 2e}\left(1-\frac{1+\chi_p(-b)}2p^{-s}\right).
\]
This series is normally convergent on $\Re s>3/2$. The products over
$p$ are bounded uniformly in $b,e$ on every compact subset of
$\Re s>1$, and $\sum_{n\ge1}r_3(n)n^{-\sigma}$ converges for
$\sigma>3/2$ by summing points of $\mathbb Z^3$ in dyadic annuli.
One has $B(2)>0$.

\begin{corollary}\label{cor:cubic-boundary}
There is a holomorphic function $H$ on $\Re s>1$ such that
\[
 Z_{\rm cub}(s)=
 \frac{\zeta(s)\zeta(s-2)\zeta(s-1)^{1/2}R(s)}{24}
 +\frac{\zeta(s)^2\zeta(s-1)B(s)}{48}+H(s).
\]
The identity holds initially for $\Re s>3$. It gives a meromorphic
continuation to $\Re s>2$ and continuation on paths in
$\Re s>3/2$ avoiding zeros and the pole of $\zeta(s-1)$, with the
usual continuation past isolated meromorphic poles. The pole at three
is simple, with residue
\[
 \rho=\frac{\zeta(3)\zeta(2)}{24}
     \prod_{p>2}\left(1-\frac{p+1}{2p^3}\right).
\]
At two, in a slit disk reached from the right,
\[
 Z_{\rm cub}(s)=\frac{\zeta(2)^2 B(2)}{48(s-2)}
       -\frac{\zeta(2)R(2)}{48\sqrt{s-2}}+O(1).
\]
In particular no meromorphic germ at two agrees with this continuation.
\end{corollary}
\begin{proof}
For the identity row, factor $\zeta(s)\zeta(s-1)\zeta(s-2)$
out of the subgroup series; its odd-prime correction is
$\prod_{p>2}(1-(p+1)p^{-s}/2)$.
Extracting $\zeta(s-1)^{-1/2}$ leaves $R$.
For every half-turn row put back the three common denominator factors
of $\zeta(s)^2\zeta(s-1)$, including at the flipped primes. The remaining
Euler factors at $p\mid e$ produce exactly $e^{-s}2^{-\omega(e)}$;
the dyadic factors give $P_b$. The quarter-, third- and sixth-turn
products are holomorphic on $\Re s>1$ by absolute convergence, since
their good factors differ from one by $O(p^{-\Re s})$ uniformly on
each compact subset. They give $H$ (including their group-order
weights). At two use $\zeta(0)=-1/2$ and
$(s-2)\zeta(s-1)=1+O(s-2)$.
\end{proof}

The decomposition also gives the unsmoothed estimate
\[
 \sum_{m\le X}a(m)=\frac{\rho}{3}X^3+
 O\big(X^2(\log(2X))^{3/2}\big).
\]
For the error term, the summatory coefficients of
$\zeta(s)\zeta(s-1)\zeta(s-2)$ are
$\zeta(3)\zeta(2)X^3/3+O(X^2\log(2X))$: sum the squares of positive
integers up to $X/(a b)$ with weight $b$ and bound the missing terms
of $\sum_{a,b}a^{-3}b^{-2}$ by $O(1/X)$.
The unsigned coefficients $h$ of the odd-prime correction satisfy
\[
 \sum_{j\le X}\frac{|h(j)|}{j^2}\le
 \prod_{2<p\le X}\big(1+(p+1)/(2p^2)\big)\ll\sqrt{\log(2X)}
\]
by Mertens' prime reciprocal estimate. Dyadic decomposition gives
\[
 \sum_{j>X}|h(j)|/j^3\ll X^{-1}\sqrt{\log(2X)}.
\]
These estimates
give the claimed error for the identity row. For the half-turn part
strip off $\zeta(s-1)$; the remaining Dirichlet expansion is absolutely
summable at two, by the same absolute bounds used for $B$ and for
$\zeta(2)^2$. Since $\sum_{j\le y}j\ll y^2$, its contribution is
$O(X^2)$. Absolute convergence at two gives the same bound for $H$.

\subsection*{An intrinsic class-number interpretation}
For any $m$, $a(m)$ is also the number of integral positive ternary
Gram classes of determinant $m^2$ with Hasse invariant $+1$ at every
odd $p\mid m$. Integrality here is with respect to the bilinear Gram matrix.
At odd $p\nmid m$, diagonalization of a unimodular matrix gives
unit entries and Hasse invariant $+1$. The real invariant is $+1$,
so Hilbert reciprocity gives the invariant $+1$ at two as well.
Hasse--Minkowski and the square determinant put the lattice in the
rational cubic space. Every full integral lattice in that space
enlarges to a self-dual lattice. At two use norm-integrality as above.
At odd $p$, a maximal integral enlargement exists inside the dual of
the original lattice. On diagonalizing this maximal lattice no exponent
can be $\ge2$, since then division of the corresponding basis vector by
$p$ gives an integral enlargement. The number of exponent-one entries
is even. If it is two the diagonal is $u,pv,pw$ with $u,v,w$ units;
its Hasse invariant is $\chi_p(-vw)$ by the odd-prime Hilbert symbol
formula. Thus the last two unit coefficients define a split residual
plane. An isotropic primitive combination of its two vectors has
norm in $p^2\mathbb Z_p$ and pairs into $p\mathbb Z_p$ with the
lattice, again permitting enlargement. Thus every maximal lattice
at odd $p$ is unimodular. Gluing gives the required unimodular
global parent and unimodular uniqueness identifies it with $L$.
The resulting index is $m$ by the determinant formula. All implications
reverse for a represented lattice.

\subsection*{Other one-class ternary genera}
The tree calculation extends to a ternary parent whose proper
genus has class number one. Let $G$ be a positive integral symmetric
ternary Gram matrix with $D=\det G$, $L$ its lattice, and
$\Gamma=\mathrm{SO}(L)$. Fix a finite set $S$ of primes containing all
divisors of $6D$. Suppose that every lattice in the proper genus of $L$
is properly isometric to $L$. Write $Z_L^{(S)}(s)$ for the subform
Dirichlet series restricted to indices prime to $S$ and
$\zeta^{(S)}(s)=\zeta(s)\prod_{p\in S}(1-p^{-s})$.
\begin{corollary}\label{cor:ternary-genus-one}
For $(m,\prod_{p\in S}p)=1$ the coefficient is the following finite sum.
Range over rational rotations $g$ with $g^2\in\Gamma$, whose least
denominator in the chosen basis is squarefree and divides $m$; call this
denominator $e$. Add $|\Gamma|^{-1}$ times the product over
$p^k\parallel m$ of the $T^k$ coefficients of
\[
 \begin{cases}
 TF_g/2 &p\mid e,\\
 F_g(1-i_p(g)T/2)&p\nmid e.
 \end{cases}
\]
Here $i_p(g)$ counts the $g$-fixed isotropic lines of $L/pL$.
The common lattice-independent stable factor $F_g$ is $F_1,F_2,
J_p(-1),J_p(-3)$ according as the rotation has order 1, 2, 4, or 3
or 6. In the second row $i_p(1)=p+1$; for a half-turn of primitive
axis $v$,
$i_p(g)=1+\chi_p(-D\,v^tGv)$, and for the remaining orders
$i_p(g)=1+\chi_p(-1)$ or $1+\chi_p(-3)$ respectively.
The series is meromorphic on $\Re s>2$ and has a simple pole at three
with positive residue
\[
 |\Gamma|^{-1}\prod_{p\in S}(1-p^{-1})\,
 \zeta^{(S)}(3)\zeta^{(S)}(2)
 \prod_{p\notin S}\big(1-(p+1)/(2p^3)\big).
\]
It has no meromorphic germ across the real point $s=2$.
\end{corollary}
\begin{proof}
At $p\notin S$ all vertices of the local self-dual tree are locally
isometric to $L_p$. Every nonsingular
ternary form over $\mathbb F_p$ contains a hyperbolic plane, and its
remaining line is determined up to squares by the determinant. A
residue-field isometry lifts to $\mathbb Z_p$: given an isometry
modulo $p^j$, a correction by $I+p^jX$ solves the next linearized
Gram equation because $X^tA+AX$ maps onto the symmetric matrices
for an invertible $A$ and odd $p$. A change of determinant sign can
be made by $-I$. Keep the parent fixed at all $p\in S$. Each
resulting global vertex is in the proper genus of $L$ and hence
is a permissible parent. For each $M$ the local tree Euler count
therefore still counts just one isometry class. The local Smith
argument identifies displacement with denominator at $p\notin S$;
at all other primes there is no displacement.

Such a $g$ has finite order because $g^2\in\Gamma$. Rationality of
the characteristic polynomial implies order $1,2,3,4$ or $6$.
The local calculations are tame as $p\nmid6$, and give the displayed
stable factors. For the half-turn the orthogonal plane has determinant
class $D/(v^tGv)$; the resulting symbol gives the stated number of
lines. For the other nonidentity orders the fixed axis is nondegenerate
and the two nontrivial eigenlines are isotropic exactly when present.
All rotations with $g^2\ne1$ form a finite list: a positive real
rotation with a prescribed nonidentity square is supported on the
same axis and has at most two possible angles. Half-turns are enumerated
by primitive integral vectors $v$ up to sign, through
$g=2v(v^tG)/(v^tGv)-I$. If $q=v^tGv$, then
$q\le 2De$: the gcd of the entries of $v(v^tG)$ equals the gcd
of the entries of $v^tG$, and multiplying by the adjugate of $G$
shows this gcd divides $D$. This gives a finite search bound for each
index. At primes dividing $e$, $g^2L=L$ and the endpoints do form
the required edge by the squarefree Smith exponents.

For analytic continuation, the finite set of rotations with $g^2\ne1$
has Euler products absolutely convergent on $\Re s>1$. For each
half-turn strip off $\zeta^{(S)}(s)^2\zeta^{(S)}(s-1)$. The remaining
term is
\[
 2^{-\omega(e)}e^{-s}
 \prod_{p\notin S,\ p\nmid e}(1-i_p(g)p^{-s}/2).
\]
Only squarefree $e$ outside $S$ occur. This series of terms is
normally convergent on $\Re s>3/2$ by $q\le2De$ and the ternary
lattice point bound. Thus all nonidentity contributions are meromorphic
in a neighborhood of two. The identity term is
\[
 |\Gamma|^{-1}\zeta^{(S)}(s)\zeta^{(S)}(s-1)\zeta^{(S)}(s-2)
       \prod_{p\notin S}(1-(p+1)p^{-s}/2).
\]
As above its product is $\zeta^{(S)}(s-1)^{-1/2}$ times a holomorphic
unit near two. Since $\zeta^{(S)}(s-2)$ has zero of order $|S|$ at
two and $\zeta^{(S)}(s-1)$ has a simple pole there, the identity
contribution is $(s-2)^{|S|-1/2}$ times a holomorphic unit, which
precludes a meromorphic germ. Evaluating at three gives the residue
as stated.
\end{proof}
\subsection*{All indices for the lattice of determinant two}
Write $Z_{112}$ for the unrestricted subform series of the lattice
with diagonal Gram matrix $(1,1,2)$. The dyadic parent change gives the following local factors.

\begin{theorem}\label{thm:det-two}
Use Theorem~\ref{thm:cubic-all} with its odd-prime rows and rotation
multiplicities unchanged. At two, include the coefficient of
$T^{v_2(m)}$ for every $m$, including odd $m$, using this table with
$F=((1-T)^2(1-2T))^{-1}$:
\[
\begin{array}{c|l}
\text{rotation}&\text{factor at two}\\ \hline
1 &(3-2T)F_1\\
2,\ b\ {\rm even}&(1+3T)F\\
2,\ b\equiv1\pmod4 &(3-T+2T^2)F\\
2,\ b\equiv3\pmod4 &(3-2T)(1+3T)F\\
4 &(1+T)/(1-T)^2\\
3,\ 6 &T/((1-T)(1-T^2)).
\end{array}
\]
The result is the coefficient of $Z_{112}$ for every $m\ge1$.
Its residue at $s=3$ is $11/4$ times the cubic residue; the same
unsmoothed error bound holds. The coefficient of $(s-2)^{-1/2}$ in
its continuation is $5/2$ times the corresponding cubic coefficient,
and the term of order $(s-2)^{-1}$ has strictly positive coefficient.
\end{theorem}
\begin{proof}
In the original lattice the vectors $(1,1,0),(1,-1,0),(0,0,1)$
give an index-two lattice $K$ of Gram matrix $2I_3$. Use its basis
and write $B=2\,\cdot$, twice the usual dot product. The local lattice
$K_2$ is intrinsic, being exactly the vectors with $B(x,x)\in2\mathbb Z_2$.
Every integral bilinear dyadic lattice $P$ of determinant valuation
one lies in $\frac12K_2$: a primitive cubic vector has norm valuation
at most one. The nonzero residues $2P\bmod 2K_2$ have weight two
since $B(x,x)$ is integral. Distinct weight-two residues have odd dot
product, so their halves cannot both belong to $P$. Covolume then
forces $K_2\subset P$ and $[P:K_2]=2$. Conversely each of the three
weight-two lines $\lambda$ gives an integral parent
$P_\lambda=K_2+\mathbb Z_2\lambda/2$, where any integral lift can be used.
All have pairwise intersection $K_2$.

Every positive integral lattice of determinant two is isometric to
the stated parent. The open ball of radius $\sqrt2$ has volume greater
than $8\sqrt2$, so Minkowski gives a unit vector. Its rank-two
orthogonal complement still has determinant two, and the disk of the
same radius has area greater than $4\sqrt2$. Splitting off another
unit gives the assertion. At odd $p$ the permissible local parents
are exactly the self-dual vertices for $B/2$. Consequently an abstract
represented lattice $M$ in this rational space can be counted by the
same contractible product of odd-prime parent trees as in the cubic
proof, with one dyadic membership test: $M_2$ must lie in at least
one $P_\lambda$. Indeed any permissible choices of odd vertices and
one such dyadic parent glue to an integral lattice of determinant
two. Conversely any embedding in the stated parent supplies these
local choices. To anchor a corner of the odd complex, glue its odd
vertices to $K_2$; the resulting lattice for $B/2$ is integral
unimodular at every prime, hence cubic. This gives exactly the same
anchoring group, rotations and odd local factors as before.

For a dyadic rotation $g$ write $F_{g,E}$ for the stable-sublattice
series in $E$. For $M_2\subset K_2$, its index in any $P_\lambda$
is twice its index in $K_2$. Every other permissible $M_2$ belongs
to just one $P_\lambda$, which must be fixed if $M_2$ is fixed.
Thus the requisite trace is
\[
 T F_{g,K_2}+
 \sum_{\lambda:\,gP_\lambda=P_\lambda}
       (F_{g,P_\lambda}-T F_{g,K_2}).
\]
This counts the identity once for every eligible lattice at this place,
regardless of how many parents contain it.

The identity fixes all three parents. A quarter-turn permutes two
coordinates modulo two and fixes just the line they span by their sum.
This enlarges its Gaussian plane by the fractional ideal
$(1-i)^{-1}\mathbb Z_2[i]$, leaving the module type unchanged.
The other rotations of orders three and six have an unramified
irreducible two-dimensional summand modulo two, whose nonzero elements
are precisely the three weight-two vectors: the fixed axis is the
nonisotropic line and no weight-two line is fixed. Indeed the weight-two
plane is the kernel of the invariant coordinate-sum form and is the
unique nontrivial residue summand.

For a half-turn of primitive axis $v$ let $E=vv^t/b$ be the positive-axis
projection. A stable ambient is split if and only if $E$ preserves it;
its two possible series are $S=(1+3T)F$ and $G=(1-2T)^{-1}+TS$
from the cubic calculation. If $b$ is even then the reduction of $g$
interchanges the two coordinates on which $v$ is odd. Its only fixed
parent adjoins $v/2$ to $K_2$. The image $E K_2$ lies in
$\mathbb Z_2v/2$, so the enlarged parent is split. For odd $b$,
$K_2$ is split, all three parents are fixed and $P_\lambda$ is split
exactly when $E\bmod2$ preserves $\lambda$. When $b\equiv1\bmod4$,
$v$ has weight one modulo two: precisely the line avoiding that
coordinate works, giving $S+2G-2TS$. When $b\equiv3\bmod4$, $v$
has weight three and is orthogonal to each weight-two residue, giving
$(3-2T)S$. These are the rows in the table.

The analytic calculation is unchanged except for these explicit
dyadic factors. The identity term acquires $3-2^{1-s}$; at $s=3$
and $s=2$ this equals $11/4$ and $5/2$. The half-turn polynomial
after removing $F$ belongs to the displayed finite list, and is
positive at $T=1/4$ in all cases. The same absolute estimates therefore
hold for its meromorphic contribution, with positive residue at two.
The finitely many remaining rotations are holomorphic near two, and
the finite dyadic change preserves the summatory error bound.
\end{proof}
For the determinant-two parent, the coefficients at indices
$5,7,35$ are respectively $9,13,234$.
\subsection*{A finite test for dyadic changes of the cubic parent}
Let $C=\mathbb Z^3$ with the cubic metric, and take any lattice $P$
with $2^aC\subset P\subset C$. A common scalar change of the metric or lattice leaves the
construction unchanged. Let $\mathcal P$
be the orbit of $P_2$ under $\mathrm{SO}(C_2)$. Suppose that the
integral signed-permutation rotation group $\Gamma$ acts transitively
on $\mathcal P$. The lattices are
given by subgroups of $C/2^aC$, and the needed residue matrices of
$\mathrm{SO}(C_2)$ can be listed by primitive quaternion coordinates
modulo $2^{a+2}$. In fact every dyadic rational rotation is given by
a quaternion by the Cayley argument above over $\mathbb Q_2$; a
primitive quaternion norm has valuation at most two. Dividing the
matrix numerators and the norm by the common power of two therefore
determines the rotation modulo $2^a$. Each primitive residue lift is
a genuine dyadic rotation. For example, modulo two the group consists
exactly of the coordinate permutations: a vector of norm one has
weight one or three, and the weight-three vector cannot be a column
of an orthonormal basis as its complement is alternating. The integral
rotations realize all permutations. The orbit hypothesis consequently
holds for \emph{every} intermediate $2C\subset P\subset C$.

\begin{corollary}\label{cor:dyadic-cubic-family}
Under the orbit hypothesis the full subform series has a finite
coefficient formula at every index, with explicit leading residue,
unsmoothed error term and real square-root branch at two.
Keep the odd local factors and the individual rational rotations in
the cubic formula, grouping half-turns only when all their local
coefficients coincide. For each rotation $g$ replace its dyadic row
by the following factor even when the index is odd:
\[
 D_g(T)=
 \sum_{\substack{\varnothing\ne J\subset\mathcal P\\gJ=J}}
 (-1)^{c_g(J)+1}\,
 T^{\log_2[P_0:E_J]}F_{g,E_J}(T).
\]
Here $E_J=\bigcap_{P'\in J}P'$, $P_0$ is any member of $J$, and
$c_g(J)$ counts the cycles of the permutation induced by $g$ on $J$.
Every factor $F_{g,E_J}$ is one of the dyadic series evaluated in
the preceding proofs. The original cubic residue at three and the
coefficient of $(s-2)^{-1/2}$ are multiplied, respectively, by
$A(1/8)$ and $A(1/4)$, where $A(T)=D_1(T)/F_1(T)$ is a polynomial
positive at both arguments. The counting error remains
$O(X^2(\log(2X))^{3/2})$.
\end{corollary}

\begin{proof}
Glue any collection of odd-prime self-dual parents to the intrinsic
dyadic cubic lattice $C_2$. This gives a global positive unimodular
lattice, necessarily cubic. In its coordinates any further choice
$P'\in\mathcal P$ corresponds to the stated parent $P$ by the assumed
transitivity of the integral rotation group. Hence the odd-prime
parent complex of a represented $M$ is unchanged, while dyadic
membership is precisely $M_2\subset P'$ for some $P'\in\mathcal P$.

If $gM_2=M_2$ the permutation $g$ preserves its set of permissible
dyadic parents. On any nonempty finite set with $t$ such cycles, the
alternating sum over nonempty unions of cycles is
$\sum_{j=1}^t(-1)^{j+1}\binom tj=1$; for an empty set it is zero.
Interchanging the resulting finite sums gives the displayed polynomial
combination of stable series. For half-turns the series is the split
or glued half-turn series according as the positive eigenprojection
preserves $E_J$ or not. For a quarter-turn project to the axis and to
the Gaussian plane. The two projected modules have ranks one over
$\mathbb Z_2$ and $\mathbb Z_2[i]$; the fiber-product classification
above gives a split ambient or a glued ambient according as the axis
projection preserves $E_J$ or not. For orders three and six there is
always a split axis and an unramified quadratic plane. Eigenprojections
can be tested on the generators of
$E_J$ by rational arithmetic and dyadic valuations.

Each half-turn row after removing $F_2$ has a polynomial drawn from
a finite list, since $\mathcal P$ is fixed. At $T=1/4$ it is positive:
the genuine local stable-lattice series converges there and contains,
for example, a sufficiently small scalar multiple of $C_2$ in every
parent. The global half-turn sum satisfies the preceding normal-convergence
estimate before axes are grouped by norm. The identity
row is $A(T)F_1(T)$. Its coefficients dominate those of $F_1(T)$
because any one of the parents by itself already contributes that
many sublattices. Hence $A(T)\ge1$ on $0<T<1/4$, and $A(1/4)\ge1$
by continuity. The singular expansion and the error estimate now follow
by the proof of Corollary~\ref{cor:cubic-boundary}.
\end{proof}
\begin{corollary}\label{cor:fcc-bcc}
The face-centered cubic and body-centered cubic lattices have identical
proper and ordinary subform Dirichlet series. They are not similar.
\end{corollary}
\begin{proof}
Use the face-centered lattice $D=\{x\in C:\sum x_i\equiv0\bmod2\}$
and the integral representative of its dual
$E=2D^*=\{x\in C:x_1\equiv x_2\equiv x_3\bmod2\}$ for the
body-centered lattice. The dual description follows because $2C\subset
D$ and its residue plane has orthogonal complement spanned by $(1,1,1)$.
Each has only one dyadic parent in its rotation orbit since every
residue rotation permutes the coordinates. At odd places their local
lattices both equal $C_p$. For each half- or quarter-turn the dyadic series
depends just on whether its rational axis projection preserves the
ambient lattice. This projection $U$ is self-adjoint for the dot
product, so $UD_2\subset D_2$ if and only if $UD_2^*\subset D_2^*$,
by the definition of the dual and reflexivity of full lattices.
Scaling by two does not change this property. The dyadic series of the
other rotation orders are independent of the stable ambient. Thus
each local factor agrees in the two computations of
Corollary~\ref{cor:dyadic-cubic-family}. There are twelve shortest
vectors in $D$ of squared length two and eight in $E$ of squared
length three; a similarity preserves the number of shortest vectors.
\end{proof}

\section{Ternary parent classes and marked neighbor trees}\label{sec:ternary-general}
Let $G$ be any positive definite symmetric rational ternary Gram
matrix, scaled to integral entries, and $\Gamma=\mathrm{SO}(L)$.
Proper and ordinary counts coincide in dimension three.
Theorem~\ref{th:incidence} gives the coefficients at every index;
at unramified odd primes its local traces can be organized on trees
whose vertices retain their global parent classes.

\paragraph{All local modules.}
Write $X_a=(1-aT)^{-1}$, $Y=(1-T^2)^{-1}$ and
$J_p(c)=X_1^3$ or $X_1Y$ according as $(c/p)=1$ or $-1$.
Let $u$ be a rotation preserving a lattice $E$. For $u\ne1$ let $U$
be the rational projection onto its positive axis. Put
$\delta=1$ if $U E_p\subset E_p$, and $\delta=0$ otherwise.
Write $S_p=(1+(p-1)T)X_1^2$ and $A_p=1+TS_p$.
The complete stable-subgroup table by order of $u$ is
\[
\begin{array}{c|l|l}
\operatorname{ord}u&\text{prime condition}&\sum_{k\ge0}b_{E,u}(p^k)T^k\\ \hline
1&\text{all}&X_1 X_p X_{p^2}\\
2&p\ne2& X_1^2 X_p\\
2&p=2,\ \delta=1&(1+3T)X_1^2 X_2\\
2&p=2,\ \delta=0&(1-T+4T^2)X_1^2 X_2\\
4&p\ne2& J_p(-1)\\
4&p=2& \delta S_2+(1-\delta)A_2\\
3&p>3& J_p(-3)\\
3&p=2&X_1Y\\
3&p=3&\delta S_3+(1-\delta)A_3\\
6&p>3&J_p(-3)\\
6&p=2&X_1Y\\
6&p=3&X_1^2 .
\end{array}
\]
Every rational proper torsion rotation in dimension three has one of
the listed orders: its quadratic complementary characteristic polynomial
is cyclotomic or a power of $X+1$. At the three exceptional
order/prime pairs $(2,2),(4,2),(3,3)$ the integral algebra generated
by $u$ is a fiber product over the residue field, and the split
condition is precisely integral preservation by $U$. The fiber-product calculation of Section~\ref{sec:cubic} applies to
an arbitrary stable $E_p$. Saturation under the two eigenprojections
gives the split module; the original module is either the split module
or a graph of residue components. In the half-turn case every graph of
a nonzero linear form on the two-dimensional negative residue plane
is conjugate under $\mathrm{GL}_2(\mathbb Z_2)$. In the Gaussian or
Eisenstein case multiplication by a unit accounts for each nonzero
residue scalar. Thus the counts are independent of any further
embedding data. The unramified and ramified direct-product rows
follow from the same Bezout splitting and the ideal series of the
degree-one or degree-two valuation rings; specifically at $(6,3)$
the factors $X-1,\Phi_6(X)$ are relatively prime over $\mathbb Z_3$.
The projections in the three exceptional cases are
$(1+u)/2$, $(1+u^2)/2$, and $(1+u+u^2)/3$, respectively. These calculations give the local formula at all primes, including
those dividing $\det G$.

Enumerate parent transporters by integral matrices $B$ satisfying
$B^tGB=m^2G$ and $\det B=m^3$, subject to the column bounds
$|B_{ij}|\le m\sqrt{G_{jj}(G^{-1})_{ii}}$, from Cauchy--Schwarz.
These matrices give the parent transporters $B/m$. Identify parents
by right multiplication by $\Gamma$ and retain only
$[L:L\cap (B/m)L]\mid m$. Each retained parent contains $mL$ and
lies in $m^{-1}L$. For a fixed torsion transporter $u$, any preserved
parent subset is a union of cycles entirely inside the finite retained
list; it must include the cycle of $L$. Successively adding or omitting
the other cycles gives the sums of Theorem~\ref{th:incidence}. The projection
test in the table uses a local lattice basis, or reduction of the
intersection modulo a sufficiently large $p$-power. 

\paragraph{Prime powers and global parent labels.}
Take $p\nmid 2\det G$. In the self-dual neighbor tree at $p$, identify
a vertex with the global lattice given by it at $p$ and $L$ at all
other places. Mark a vertex if this global lattice is properly
isometric to $L$. For a finite nonempty subtree $\mathcal T$, write $\partial\mathcal T$ for its set
of leaves (including the unique vertex of a singleton), and let
$E_{\mathcal T}$ be the global intersection of its vertices. Write
$p^{j_{\mathcal T}}=[L:E_{\mathcal T}]$ for trees containing $L$.
\begin{theorem}\label{th:ternary-hull}
For every $k\ge0$ the actual coefficient at $p^k$ is
\[
 \frac1{|\Gamma|}
 \sum_{\mathcal T}\frac1{|\partial\mathcal T|}
  \sum_{u\in\mathrm{Stab}^+(\mathcal T)}
      (-1)^{c(u,\partial\mathcal T)-1}
      b_{E_{\mathcal T},u}\big(p^{\,k-j_{\mathcal T}}\big).
\]
Sum over subtrees with $L$ a leaf, $j_{\mathcal T}\le k$, every leaf marked, and
no marked interior vertex of the chosen subtree. The stabilizer consists of global
proper rational isometries preserving $\mathcal T$ and fixing $L_\ell$
at all places $\ell\ne p$; $c$ counts cycles on leaves. This is a
finite sum, with $b$ given by the table above. If the genus has
one class, only singletons and edges survive.
\end{theorem}
\begin{proof}
In the unanchored incidence sum proved in Theorem~\ref{th:incidence},
take the convex hull of each parent set $F$ in the $p$-tree. Its
leaves must lie in $F$. Conversely every subset of its marked
vertices containing its leaves spans the tree. Intersecting all the
vertices of the hull gives exactly the intersection of the original
subset, since the unique geodesic between two vertices contains
their intersection in every one of its lattices. Every hull vertex is
on such a geodesic or is itself a leaf.

Within a fixed hull, orbit--stabilizer rewrites the unanchored
sum over orbits of $F$,
$\sum_{[F]}|\mathrm{Stab}(F)|^{-1}\sum_{u\in\mathrm{Stab}(F)}$
as $|\mathrm{Stab}(\mathcal T)|^{-1}\sum_{u\in\mathrm{Stab}(\mathcal T)}
\sum_{F:uF=F}$, with the same integrands. For each $u$ the optional
marked interior vertices form a union of cycles. Including an
optional cycle changes the sign; the sum therefore vanishes if any
such vertex exists. Otherwise only the leaf set remains. Marking a
leaf as $L$ replaces the unanchored stabilizer denominator by
$|\partial\mathcal T|\,|\Gamma|$ by the same orbit-stabilizer identity. The
restriction $j_{\mathcal T}\le k$ bounds every vertex's distance to $L$ by
$k$, by the symmetric Smith indices of the self-dual tree, so the
tree sum is finite.
\end{proof}

The leaf sets also admit a graph-theoretic description. In the ball of
radius $k$ around $L$, join two marked vertices $P,Q$ in an auxiliary
graph if there is no third marked vertex $R$ satisfying
$\mathrm{dist}(P,R)+\mathrm{dist}(R,Q)=\mathrm{dist}(P,Q)$.
The leaf sets in the formula are exactly the cliques containing $L$
whose intersection index divides $p^k$. Indeed the ball is convex,
every interior vertex lies on a path between two leaves, and any
member of a proposed leaf set that were interior would lie between
two other members. The distances are obtained from the denominator of
the relative transporter in a basis of $L$ by the self-dual Smith
exponents. The marked vertices therefore determine the leaf sets.

At $k=1$ the formula reduces to a class-sensitive expression. Let
$A_\Gamma(p)$ count the $\Gamma$-orbits of hyperplanes of $L/pL$.
Let $t_p$ be the number of double cosets $\Gamma g\Gamma$ of
rational proper isometries with $[L:L\cap gL]=p$, and $r_p$ the
number of these cosets fixed by inversion. Then
\[
                a_L(p)=A_\Gamma(p)-\tfrac12(t_p-r_p).
\]
Indeed the double cosets are the orbits of marked outgoing neighbor
edges. Reversal induces inversion of double cosets. The intersection
child of a fixed orbit contributes one ambient orbit and one true
class. Each pair of distinct reversed orbits contributes two ambient
orbits and one true class. All other hyperplane children have just
one marked parent. Even two distinct incident neighbors have distinct
intersection hyperplanes by the local tree construction. For any
$p$ here, including $3$, the first term is the average over $u\in\Gamma$
of $\sum_{\alpha\in\mathbb F_p^\times}
(p^{\dim\ker(u-\alpha)}-1)/(p-1)$, since it counts invariant
hyperplanes via eigenlines in the dual. For $p>3$ the contributions
are $p^2+p+1$, $p+2$, $2+(-1/p)$ and $2+(-3/p)$ for orders
$1$, $2$, $4$ and $3$ or $6$, respectively.

For indices supported on several good primes the same cancellation
has an exact finite form. Exclude $m=1$, and for each $p\mid m$ take a
finite subtree $\mathcal T_p$. Put $\mathcal Q=\prod_{p\mid m}\mathcal T_p$, with $L$ prescribed
away from these primes. Mark a tuple if its global lattice is properly
isometric to $L$. For $u$ preserving the box and the external local
lattices the unanchored weight is
\[
 \sum_{\mathcal U}(-1)^{|\mathcal U|-1}.
\]
Here $\mathcal U$ runs over collections of cycles of $u$ on marked
tuples, whose projections together contain every leaf of every factor
(a singleton counts as a leaf). The whole contribution is this weight
times $b_{E_{\mathcal Q},u}(m/[P:E_{\mathcal Q}])$, summed over $u$ and divided by the
global proper stabilizer order of $\mathcal Q$; $P$ is any vertex and $E_{\mathcal Q}$ the
intersection of all vertices. Use one representative of each global
box class that contains a marked tuple, with $[P:E_{\mathcal Q}]\mid m$.
Indeed coordinatewise convex hull gives the same intersection at
each prime; a union of cycles spans the box precisely by this leaf
test. Marking one tuple bounds all relevant trees by the corresponding
index exponents and hence proves finiteness, as above. The parent
class marks occur on the global tuples, and need not factor as
independent choices in the local trees. At the remaining primes the
all-index incidence expression and the explicit stable-module table
apply without any tree assumption.

For $L=\langle1,1,11\rangle$ and $p=3$, the line spanned by $(1,0,1)$
modulo three is isotropic. The corresponding neighbor cannot be
isometric to $L$: its intersection hyperplane omits $e_1$, and the
only rational unit vectors with denominator dividing three are
$\pm e_1,\pm e_2$ (multiply the norm equation by nine). Thus it
has at most two unit vectors whereas $L$ has four. The local
neighbor is nevertheless in the proper genus. Consequently it is an
unmarked vertex in the theorem.

For $G=\operatorname{diag}(1,1,7)$ and $p=3$, a marked vertex
can occur beyond an unmarked intermediate vertex. Every rational
unit vector with denominator dividing three is one of $\pm e_1,\pm e_2$ or
$(\pm1,\pm1,\pm1)/3$. No vector of the latter type is perpendicular
to either unit coordinate vector or any other vector of the latter
type. Thus a proper isometry of denominator dividing three must send
$e_1,e_2$ to the coordinate vectors and hence be integral. No adjacent
vertex is marked. But
\[
 A=\begin{pmatrix}-7&-2&-14\\-2&-7&14\\-2&2&5\end{pmatrix},
 \qquad A^tGA=81G,\qquad\det A=729 .
\]
The rational rotation $A/9$ has exact denominator nine, so its marked
vertex lies at distance two. The path between these two marked
endpoints has precisely an unmarked interior vertex.

\section{Orientation in even rank}\label{sec:quaternary}
In this section write $a_L(m)$ and $a_L^+(m)$ for the ordinary
and proper coefficients, respectively. A chiral ordinary class has
two proper orientations in the rational space; the parent may
represent just one.

\subsection{Fixed sublattices of a sign change}
Put $b_t(m)=\#\{M\subseteq\mathbb Z^t:[\mathbb Z^t:M]=m\}$, with $b_0(1)=1$ and $b_0(m)=0$ for $m>1$. Write $*$ for Dirichlet convolution.

\begin{lemma}\label{lem:split-sign-glue}
Let $F_{r,t}(m)$ count the index-$m$ lattices fixed by $(x,y)\mapsto(x,-y)$ on $\mathbb Z^r\oplus\mathbb Z^t$. Then
\[
 F_{r,t}(m)=
 \sum_{\substack{0\le j\le \min(r,t)\\2^j\mid m}}
 \genfrac{[}{]}{0pt}{}{r}{j}_{2}\genfrac{[}{]}{0pt}{}{t}{j}_{2}\,
 |\operatorname{GL}_j(\mathbb F_2)|\,(b_r*b_t)(m/2^j),
\]
where the Gaussian binomial counts codimension-$j$ subspaces over $\mathbb F_2$ and the group of rank zero has order one.
\end{lemma}
\begin{proof}
For a full lattice $M$ let $R_+,R_-$ be its projections and $K_+,K_-$ its intersections with the two summands. Lifting $x\in R_+$ to $(x,y)\in M$ defines an isomorphism $R_+/K_+\longrightarrow R_-/K_-$ by $x\mapsto y$. Its graph recovers $M$ and each graph gives one such lattice. The sign change preserves this graph precisely when negation acts trivially on the common quotient. This makes the quotient an $\mathbb F_2$-space, of some dimension $j$. For fixed projections the choices of kernels correspond bijectively to codimension-$j$ subspaces of $R_\pm/2R_\pm$, followed by a choice of isomorphism. The graph has index $2^j$ in $R_+\oplus R_-$. Summing over projection indices gives the result.
\end{proof}
The Dirichlet series of $F_{r,t}$ thus factors as the product
of those of $b_r$ and $b_t$, multiplied by the displayed polynomial
in $2^{-s}$ describing dyadic gluing.

\subsection{A quaternary local extension criterion}
\begin{lemma}\label{lem:quat-integral-norm}
Let $D>1$ be squarefree and let $q_1,q_2$ be positive definite integral binary forms. On $L=\mathbb Z^2\oplus\mathbb Z^2$ take $Q(x,y)=q_1(x)+Dq_2(y)$. Let $P$ consist of primes $p\mid D$ where both binary forms are anisotropic modulo $p$. At every index $m$ supported on $P$, every rational isometry carrying one index-$m$ sublattice of $L$ onto another preserves $L$.
\end{lemma}
\begin{proof}
For nonzero $x\in\mathbb Q_p^2$, write $k=\min_i v_p(x_i)$. Anisotropy gives $v_p(q_i(x))=2k$. Nonzero terms in $q_1(x)+Dq_2(y)$ have valuations of different parity. Their sum is integral exactly when both $x,y$ lie in $\mathbb Z_p^2$, also when one input vanishes. Thus $L_p$ is exactly the set of vectors of integral norm, including at $p=2$ when the hypothesis holds. All rational isometries preserve $L_p$ at these primes. At primes outside the support of $m$, both sublattices localize to $L_p$. Equality at every finite prime gives equality of the integral lattices.
\end{proof}

\begin{theorem}\label{thm:quat-sign-family}
Fix a nonempty finite set $P$ of odd primes and put $D=\prod_{p\in P}p$. Choose $b>0$ prime to $D$ such that $1,b,D,Db$ are distinct and $-b$ is a nonsquare at every $p\in P$. Such $b$ exist by the Chinese remainder theorem. Let $L$ have diagonal form $\langle1,b,D,Db\rangle$. At every index $m$ supported on $P$,
\[
 a_L(m)=\tfrac18(b_4+4F_{1,3}+3F_{2,2})(m),\qquad
 a_L^+(m)=\tfrac14(b_4+3F_{2,2})(m).
\]
At $p\in P$ one has $a_L^+(p)-a_L(p)=((p-1)/2)^3$. At every $p^e$, $e\ge1$, the difference is at least $(\varphi(p^e)/2)^3$ and
\[
 a_L(p^e)\sim b_4(p^e)/8,\qquad a_L^+(p^e)\sim b_4(p^e)/4\qquad (e\to\infty).
\]
For distinct $p,q\in P$ both coefficients satisfy $a(pq)>a(p)a(q)$.
\end{theorem}
\begin{proof}
In a diagonal lattice with distinct positive weights, the shortest nonzero vectors are the two signed vectors on the lightest axis. Restriction to the perpendicular complement and induction give the entire integral isometry group, the independent sign changes. Lemma~\ref{lem:quat-integral-norm} applies with $q_i=x^2+by^2$, so the required coefficients are exactly the orbit counts under this sign group and its determinant-one subgroup. The full group contains two central signs, eight changes negating one or three coordinates, and six changes negating two. Burnside's formula and Lemma~\ref{lem:split-sign-glue} give the displayed coefficients.

Kernels of surjections $L\to\mathbb Z/p^e$ are parametrized by primitive rows modulo unit scaling. Rows whose four entries are units give $\varphi(p^e)^3$ distinct kernels. A sign fixing such a row projectively has all signs equal. Its stabilizer lies in the proper group and the ordinary orbit, of size eight, splits in two. At $e=1$ every other row has a zero coordinate whose sign change provides an improper stabilizer, proving exactness there. The Euler factor for $b_j$ at $p$ is $\prod_{i=0}^{j-1}(1-p^i X)^{-1}$ by column Hermite normal form. Partial fractions give $b_4(p^e)\sim p^{3e}\prod_{i=1}^3(1-p^{-i})^{-1}$. At odd $p$, the same product formula gives $F_{1,3}(p^e)=O_p(p^{2e})$ and $F_{2,2}(p^e)=O_p((e+1)p^e)$.

For any fixed sign $g$, the count $f_g$ of its fixed lattices is multiplicative, by the prime-wise decomposition of a finite-index sublattice. At a prime $p\in P$, the three sign types above have counts $p^3+p^2+p+1$, $p^2+p+2$, and $2p+2$, in strictly decreasing order. This ordering holds for $q$ as well. Each of the two groups has at least two of these types. The difference between the average of $f_g(p)f_g(q)$ and the product of the averages is
\[
 \frac{1}{2|G|^2}\sum_{g,h\in G}
 (f_g(p)-f_h(p))(f_g(q)-f_h(q))>0.
\]
\end{proof}

\begin{theorem}\label{thm:quat-missing-mate}
Put $B(x,y)=3x^2+xy+5y^2$. Let $D>1$ be squarefree with every prime divisor either $2$ or an odd prime at which $-59$ is a nonsquare. For $L=(\mathbb Z^4,B\perp D B)$ and every $m$ supported on the prime divisors of $D$,
\[
 a_L(m)=a_L^+(m)=\tfrac12\bigl(b_4(m)+F_{2,2}(m)\bigr).
\]
Every ordinary class counted here admits no improper automorphism, and exactly one of its two proper classes is represented inside $L$.
\end{theorem}
\begin{proof}
Completing the square gives $B=3(x+y/6)^2+59y^2/12$. Its shortest vectors are $\pm e_1$. Any integral isometry has columns $(\epsilon,0),(k,\delta)$ where $\epsilon,\delta=\pm1$. Comparing the cross coefficient gives $6k=\epsilon-\delta$, hence $k=0,\delta=\epsilon$. In $L$, vectors of norm at most five span the first rational plane. The other plane has minimum $3D>5$. Thus every symmetry preserves both planes, and its two restrictions are independent scalar signs. In particular every symmetry is proper. The odd-prime hypotheses ensure anisotropy of $B$ by its discriminant, and modulo two one has $B=x^2+xy+y^2$, also anisotropic. Lemma~\ref{lem:quat-integral-norm} makes the integral symmetry group control all the relevant equivalences, including all automorphisms of each sublattice. The central pair fixes every lattice; the two remaining elements each have two positive and two negative axes. Burnside averaging gives the displayed formulas. An improper rational reflection exists in the ambient space. If its image of a proper class were represented in $L$, composition with the proper equivalence would give an improper isometry between two of our sublattices, contradicting the extension lemma.
\end{proof}
For $D=2$ the coefficients are $(a_L(2),a_L(4))=(15,119)$; for $\langle1,4,3,12\rangle$ they satisfy $(a_L(3),a_L^+(3))=(15,16)$.

\subsection{A fixed diagonal lattice in every even rank}
Let $L$ now have distinct positive diagonal weights $w_1,\ldots,w_n$ with $n$ even. Let $p$ be an odd prime not dividing $\prod_i w_i$ and $e\ge1$. A cyclic index-$p^e$ sublattice is the kernel of a primitive covector modulo $p^e$ up to unit scaling. Call this kernel nonsingular if the dual quadratic norm of the covector is nonzero modulo $p$.

\begin{lemma}\label{lem:quat-cyclic-recovery}
For any lattice unimodular at an odd prime $p$, a rational isometry between an index-$p^e$ nonsingular cyclic sublattice and any index-$p^e$ sublattice of the same parent preserves that parent.
\end{lemma}
\begin{proof}
Use the symmetric pairing at $p$ to lift the covector to a vector $v\in L_p$ of unit norm. Then $L_p=H\perp\mathbb Z_p v$ with $H$ unimodular, and the kernel is $M_p=H\perp p^e\mathbb Z_p v$. Hence $M_p^\#/M_p$ is cyclic of order $p^{2e}$. The parent is an integral overlattice, so its quotient by $M_p$ is the unique subgroup of order $p^e$ in this cyclic group. The same intrinsic recovery holds on the target side since its discriminant module is isomorphic and its given parent is integral. At other primes both lattices localize to the parent. Thus rational isometry preserves the nonsingular stratum.
\end{proof}

For $S\subseteq\{1,\ldots,n\}$, set $A_\varnothing=0$. For nonempty $S$ of size $h$ put
\[
 A_S=p^{(e-1)(h-1)}
 \left[p^{h-1}-
 \begin{cases}
 \displaystyle\left(\frac{(-1)^{h/2}\prod_{i\in S} w_i}{p}\right)p^{h/2-1},& h\ \mathrm{even},\\
 0,&h\ \mathrm{odd}.
 \end{cases}\right].
\]
\begin{theorem}\label{thm:even-good-gap}
The contribution of the nonsingular cyclic stratum to $a_L^+(p^e)-a_L(p^e)$ is exactly
\[
 2^{1-n}\sum_{S\subseteq\{1,\ldots,n\}}(-1)^{n-|S|} A_S.
\]
It is a lower bound for the whole difference and is strictly positive for $p\ge5$. For fixed $L$, as $p\to\infty$ through primes,
\[
 a_L(p)=2^{1-n}p^{n-1}+O_n(p^{n-2}),\qquad
 a_L^+(p)=2^{2-n}p^{n-1}+O_n(p^{n-2}).
\]
\end{theorem}
\begin{proof}
By Lemma~\ref{lem:quat-cyclic-recovery} the ordinary and proper actions on this stratum are precisely those of the sign group and its even subgroup. A primitive row has at least one unit coordinate, which forces its projective stabilizer scalar to be $\pm1$. At each nonzero coordinate modulo $p^e$ its sign must equal this scalar since $2$ is a unit. Thus an ordinary orbit splits if and only if every coordinate is nonzero modulo $p^e$. In that case its size is $2^{n-1}$.

The quantity $A_S$ counts nonsingular projective covectors whose
coordinates vanish outside $S$. Their $h$-variable dual quadratic form
has determinant squareclass $(\prod_{i\in S}w_i)^{-1}$, whose
quadratic character is the same as that of $\prod_{i\in S}w_i$.
The nonisotropic-line formula in Section~\ref{sec:core}, applied in
dimension $h$, therefore gives the bracketed factor in $A_S$.
Each such line has $p^{(e-1)(h-1)}$ projective lifts, obtained by
normalizing one unit coordinate. Thus the displayed expression for
$A_S$ incorporates all lifts to modulus $p^e$. Inclusion--exclusion now imposes nonzero entries modulo $p^e$ at every coordinate. The forgetful map shows that all remaining ordinary classes
contribute nonnegatively to the difference. If $p\ge5$, set all but one coordinate to one and vary the last among units modulo $p$. Its square takes at least two different values, so at least one choice has nonzero norm. Every lift with all coordinates units still qualifies. Finally at index $p$ the excluded coordinate hyperplanes and the isotropic quadric contain $O_n(p^{n-2})$ projective points by the same count. The remaining $p^{n-1}+O_n(p^{n-2})$ points all have scalar stabilizer and lie in saturated orbits, proving both asymptotics.
\end{proof}

For each even rank $2k$, take $\mathop{\perp}\limits_{j=0}^{k-1}2^j B$,
where $B$ is the preceding binary form. In the complement of the first $j$ planes the vectors of norm at most $5\cdot 2^j$ span the next plane. Induction proves that its full symmetry group consists of scalar signs on the binary blocks and hence is proper. At every odd prime of unimodularity, Lemma~\ref{lem:quat-cyclic-recovery} consequently gives one-sided orientation and no improper automorphisms on the entire nonsingular cyclic stratum, at all exponents. This stratum is nonempty, since a nondegenerate quadratic form in odd characteristic cannot vanish identically.

\subsection{Orientation across several primes}
For the same diagonal $L$, let $m$ be any positive integer prime to $2\prod_i w_i$. Consider the saturated stratum of lattices whose component at each $p^e\Vert m$ is cyclic and nonsingular. For a nonempty set $S$ put
\[
 N_{p,e}(S)=\sum_{T\subseteq S}(-1)^{|S|-|T|}A_T(p,e).
\]
Here $A_T(p,e)$ is the preceding finite-field expression. Given a family $\mathcal S=(S_p)_{p\mid m}$, make a graph on $\{1,\dots,n\}$ by joining any two distinct vertices occurring in a common $S_p$. Include isolated vertices and write $c(\mathcal S)$ for the number of connected components. Define $E(\mathcal S)=1$ if every component has even size, and $E(\mathcal S)=0$ otherwise.

\begin{theorem}\label{thm:even-orientation-crt}
Let $a_{\rm ns}$ and $a_{\rm ns}^+$ count ordinary and proper classes in this stratum. Then
\begin{align*}
 a_{\rm ns}(m)
 &=\sum_{\mathcal S}2^{c(\mathcal S)-n}
            \prod_{p^e\Vert m}N_{p,e}(S_p),\\
 a_{\rm ns}^+(m)-a_{\rm ns}(m)
 &=\sum_{\mathcal S} E(\mathcal S)\,2^{c(\mathcal S)-n}
            \prod_{p^e\Vert m}N_{p,e}(S_p).
\end{align*}
The sums run over all assignments of nonempty $S_p$. In rank four there are pairs of local kernels, each admitting an improper stabilizer, whose combined global kernel has two represented orientations and no improper stabilizer.
\end{theorem}
\begin{proof}
The Chinese remainder theorem identifies these lattices with independent choices of primitive projective rows at the primes dividing $m$. Primary decomposition recovers each local congruence kernel from
their intersection. Recovery of the unimodular parent in Lemma~\ref{lem:quat-cyclic-recovery} holds at each of these primes, and outside their support the parent is unchanged. Therefore distinct global tuples are equivalent exactly under an element of the global sign group, with the determinant condition for proper equivalence.

The number of tuples with local coordinate supports exactly $\mathcal S$ is the displayed product of $N_{p,e}(S_p)$ by inclusion--exclusion. At a given prime the scalar of a stabilizing sign, found from a unit entry, forces equal signs at all nonzero coordinates modulo $p^e$. This holds simultaneously at every prime precisely when the sign is constant on each component of the graph. The stabilizer has cardinality $2^{c(\mathcal S)}$, and contains an improper element precisely when one component has odd size. The first formula weights each tuple by the inverse size of its orbit. Exactly the indicated orbits split under the determinant-one subgroup, giving the second formula. For the last assertion take different primes $p,q\ge 5$ of unimodularity and supports $S_p=\{1,2,3\}$, $S_q=\{3,4\}$. Nonsingular rows with these exact supports exist since one can vary one nonzero square while fixing the other entries. Each support misses a coordinate, while the two together connect all four vertices.
\end{proof}

\begin{corollary}\label{cor:even-separated-window}
For any integer $M\ge1$ choose positive weights $w_{j+1}>M^2w_j$. With $L=\langle w_1,\ldots,w_n\rangle$, every index $m\le M$ satisfies
\[
 a_L(m)=2^{-n}\sum_{h=0}^n\binom nh F_{n-h,h}(m),\qquad
 a_L^+(m)=2^{1-n}\sum_{\substack{0\le h\le n\\2\mid h}}\binom nh F_{n-h,h}(m).
\]
\end{corollary}
\begin{proof}
If $T$ carries one index-$m$ sublattice to another, then $mTe_j\in L$ and has norm $m^2w_j<w_{j+1}$ for $j<n$. Its last $n-j$ entries therefore vanish. Thus $T$ stabilizes the coordinate flag, and orthogonality forces $T$ to be diagonal with entries $\pm1$. All these signs preserve the parent. Apply Burnside's formula and Lemma~\ref{lem:split-sign-glue}.
\end{proof}


\section{Prime-index sublattices in varying rank}
Let \(a_p(n)\) be the intrinsic ordinary coefficient at index \(p\) of the standard lattice \(\mathbb Z^n\), and let \(b_p(n)\) be the number of signed-permutation orbits of the same sublattices. Set both sequences equal to zero at \(n=0\).
\begin{proposition}
For any odd prime \(p\) and all \(n\ge p^2\),
\[
 a_p(n)=b_p(n)-b_p(p^2)+a_p(p^2).
\]
Writing \(k=(p-1)/2\), one has
\[
 \sum_{n\ge0} b_p(n)z^n=\frac1{1-z}
 \left(\frac1k\sum_{d\mid k}\frac{\varphi(d)}{(1-z^d)^{k/d}}-1\right).
\]
Thus \(a_p(n)\) is eventually a quasipolynomial of degree \(k\) and period dividing \(k\). The proper count equals the ordinary count when \(n>k\). For \(p=2,3\) both counts equal \(n\) in every positive rank, and for \(p=5,n\ge3\) they equal \(n+\lfloor n^2/4\rfloor\).
\end{proposition}
\begin{proof}
Index-\(p\) kernels correspond to projective normal vectors. If \(w\) coordinates of a normal are nonzero, the norm-one vectors of its kernel span exactly the remaining \(n-w\) coordinate axes. Their orthogonal complement is the active kernel in \(\mathbb Z^w\), which has no norm-one vector. Hence \(w\) is intrinsic, and adjoining inactive axes preserves
the distinction between active kernels. An isometry between two active kernels has matrix \(A/p\) with integral \(A\), since \(p\mathbb Z^w\) lies in the source. For a nonintegral isometry some row of \(A\bmod p\) is a nonzero multiple of its source normal. This row has \(w\) nonzero integer entries and squared norm \(p^2\), so \(w\le p^2\). Hence increasing the rank beyond \(p^2\) introduces no new identifications. For \(n>k\) there is an inactive axis or a pair of coordinates equivalent up to sign. A sign change of the former or a signed transposition of the latter is an improper automorphism fixing the kernel. The same conclusion is immediate at \(p=2\).

Modulo signed permutations, a nonzero normal is specified by multiplicities of the \(k\) pairs of nonzero residues up to sign. Projective rescaling rotates these types by the cyclic group of order \(k\). Its elements of order \(d\) have \(k/d\) cycles of length \(d\), giving the displayed series after removing the zero tuple and adding inactive axes. The identity has pole order \(k+1\) at one and all other summands have lower pole order. All poles are \(k\)-th roots of unity. For \(p=2,3\) one active type remains and the norm-one count already distinguishes all possibilities. At \(p=5\), write \(z,u,v\) for the counts of zero coefficients and the two nonzero types. There are \(2z\) vectors of norm one and \(4\binom z2+2(\binom u2+\binom v2)\) of norm two. These recover the unordered pair \(\{u,v\}\) and hence the orbit. Summing \(\lfloor w/2\rfloor+1\) over \(1\le w\le n\) gives the last formula.
\end{proof}


\section{Indefinite spaces of rank at least three}\label{sec:indef}
Let $V$ have indefinite signature and rank $n\ge3$. The HNF
bound and the orientation-union identity continue to hold, while
integral stabilizers may be infinite. We use local isometry types
and embedding spinor norms to count the represented classes. Set
\[
 c_p(e)=\#\{\text{local types of index-}p^e\text{ sublattices of }L_p\},\quad
 c(m)=\prod_{p^e\Vert m}c_p(e),\quad c_p(0)=1 .
\]

\subsection{The spinor norm and represented classes}
The adelic spinor-genus description and Eichler's theorem give
one proper class in each proper spinor genus of an indefinite
rational lattice of rank at least three \cite[\S9.2]{SP}. More explicitly, write $\theta$ for the spinor norm, normalized on a reflection in $v$ by $Q(v)$ modulo squares. On the proper group this convention is independent of a scalar in the quadratic form. For an integral lattice $K$ in this space the proper classes in its genus are parametrized, with $K$ as origin, by the finite elementary abelian group
\begin{equation}\label{eq:spinC}
 C(K)=\mathbb A_f^\times/
 \bigl(\mathbb Q^\times\mathbb A_f^{\times2}
       \prod_p\theta(SO(K_p))\bigr).
\end{equation}
Products of local subgroups are taken in the ideles, with
squareclasses lifted to multiplicative groups. There is no positivity
restriction at the indefinite real place. Strong approximation for the spin group and the spinor-norm
image statements in \cite[\S9.2]{SP} apply to the indefinite space.

The ambient spinor-norm image follows from an indefinite ternary
subspace. Its even Clifford algebra is a quaternion algebra
split at infinity, and conjugation on the ternary space has spinor norm
the reduced norm modulo squares. Every nonzero rational number is a
reduced norm, and the reduced norm is surjective at every finite place.

Improper local isometries introduce no additional genera: choose a
vector $v$ of minimal valuation of nonzero $Q(v)$ in the integral local
lattice. Then $Q(v)$ divides $Q(v+w)-Q(v)-Q(w)$ for every $w$, so
reflection in $v$ is an integral improper isometry, also at two.

For each prime dividing $m$, choose a local type and an embedding
into $L_p$, taking $L_p$ at all other primes. Chinese-remainder
patching gives a reference lattice $K\subset L$ of index $m$.
Different tuples have distinct genera, and all $c(m)$ tuples occur. In \eqref{eq:spinC} let $T(K,L)$ be the set of products of the images under $\theta$ of the sets
\[
 \{g_p\in SO(V_p):g_p K_p\subset L_p\}.
\]
Use integral elements at almost all places. The represented proper classes in this genus are counted by $|T(K,L)|$: patching the local images yields the corresponding spinor class, and every represented class arises from such local embeddings. 
One $g_p$ per local image suffices, since the stabilizer has
trivial image in $C(K)$. For $p\nmid m$, the local image equals $L_p$ and the stabilizer
has trivial image in $C(K)$.

For the ordinary count fix a rational reflection $r$ and choose local improper isometries $j_p\in O(K_p)$, equal to $r$ wherever $r$ is integral except at finitely many places. Put $\delta_K=[\theta((rj_p)_p)]$. This class is independent of the choices of $j_p$; changing $r$ changes the product by a rational proper factor. Reversal acts on $C(K)$ by translation by $\delta_K$, since conjugation by an orthogonal transformation preserves the spinor norm. Thus the number of represented ordinary classes is the cardinality of the image of $T(K,L)$ in $C(K)/\langle\delta_K\rangle$. Summing these formulas over the $c(m)$ tuples gives both actual coefficients at all indices.

\begin{lemma}\label{lem:indlocal}
For $p\nmid 2\det B_L$, $c_p(1)=3$. Each of these types forms a single $SO(L_p)$-orbit of sublattices. For each child $K_p$ its full integral group $SO(K_p)$ has all unit squareclasses as spinor norms. For the isotropic hyperplane type $SO(K_p)$ has all local spinor norms. Moreover $\theta(SO(L_p))$ consists of unit squareclasses.
\end{lemma}
\begin{proof}
A kernel is a hyperplane modulo $p$. Normals are isotropic, of square unit norm, or of nonsquare unit norm. The nonsingular finite quadratic space of rank at least three contains a hyperbolic plane and represents both unit classes. Witt extension gives transitivity on each set of normal lines under the finite orthogonal group. Isometries over $\mathbb F_p$ are products of reflections in nonisotropic vectors; lift each vector to a vector with unit norm to obtain an integral isometry. The determinant can be corrected in a hyperplane stabilizer: in the nonisotropic case use a unit vector in the perpendicular space, and in the isotropic case use the complement of a hyperbolic plane, of rank $n-2\ge1$. Thus each set is a single proper parent orbit.
At a nonisotropic line the kernel has a unimodular rank $n-1$ summand and a unary summand scaled by $p^2$. The latter's unit determinant squareclass distinguishes the two choices, by its intrinsic top Jordan constituent. 

At an isotropic line lift a hyperbolic basis by Hensel lifting and orthogonal splitting. Its kernel replaces the hyperbolic plane $\mathbb H_p$ by $p\mathbb H_p$ (scale of its form) and keeps a unimodular complement. This has a different Jordan decomposition. 

Any unimodular summand of rank at least two represents all unit squareclasses: over $\mathbb F_p$ this follows for a split binary plane by multiplication and for an anisotropic binary plane by the finite-field norm. Lifts preserve the unit squareclass. Products of unit-norm reflections then give every unit spinor class. The same argument in $p\mathbb H_p$ gives all units by products of reflections of norms $pu,pv$; these are integral on the scaled plane. Composing one such reflection with a unit-norm reflection in the complement adds an odd valuation class. 

Finally at a unimodular odd-prime lattice every integral orthogonal transformation is a product of unit-norm reflections. For a unit diagonal basis vector $v$ at least one of $hv\pm v$ has unit norm, since the sum of their norms is $4Q(v)$. Use that reflection, if necessary the reflection in $v$, then induct on its unimodular complement. This proves the last assertion.
\end{proof}
Consequently, for every indefinite parent of rank $n\ge3$,
\begin{equation}\label{eq:indsquarefree}
 a_O(t)=a_{SO}(t)=3^{\omega(t)}\qquad(t\ {\rm squarefree},\ (t,2\det B_L)=1).
\end{equation}
For a fixed tuple, at $p\mid t$ the transporter from the reference child is $SO(L_p)SO(K_p)$ by the single-orbit assertion. Its spinor norm has trivial image in $C(K)$. Elsewhere the same conclusion holds trivially. Thus exactly one proper class is represented per tuple.

\subsection{The leading constant in indefinite signature}
\begin{theorem}\label{th:indmean}
For each of $H=O,SO$, the abscissa is $1$ and
\[
 \sum_{m\le X}a_H(m)\sim \frac{C_L}{2}X(\log X)^2,\qquad
 C_L=\prod_p\left[(1-p^{-1})^3\sum_{e\ge0}\frac{c_p(e)}{p^e}\right]>0,
\]
and in particular $(s-1)^3Z_H(s)\to C_L$ as real $s\downarrow1$.
\end{theorem}
\begin{proof}
Let $S$ contain the primes of $2\det B_L$. Local Jordan splitting bounds $c_p(e)\ll_{n,L}(1+e)^n$ uniformly: the determinant valuation is $v_p(\det B_L)+2e$, and there are at most $(1+v_p(\det B_L)+2e)^n$ ordered scale lists counted with repetition for the $n$ basis positions. At odd $p$ a unimodular constituent is specified by its rank and its unit determinant squareclass. At two there are boundedly many unimodular types of bounded rank, by the dyadic classification \cite{AGM}. These bounds include imprimitive children. 

The multiplicative $c(m)$ therefore has series $\zeta(s)^3 V(s)$ with $V(s)$ an absolutely convergent Dirichlet series for $\Re s>1/2$: at primes outside $S$ the linear Euler coefficient cancels by Lemma~\ref{lem:indlocal}. $V(1)=C_L$ is positive by absolute convergence and positive local values. The elementary estimate $\sum_{j\le x}d_3(j)=x(\log x)^2/2+O(x(1+\log x))$ follows by summing $\lfloor x/(ab)\rfloor$ and the harmonic sum twice. Convolution and dominated convergence give $\sum_{m\le X}c(m)\sim C_L X(\log X)^2/2$.

Write $m=dt$ where $d$ contains all powers at $S$ and all prime powers of exponent at least two outside $S$; $t$ is squarefree and coprime to $S_d=S\cup\{p:p\mid d\}$. At every place outside $S_d$ the local spinor norm of a child stabilizer contains all units. A finite idele may be divided by a positive rational number to make it a unit at all finite places. Its remaining class is generated by the unit squareclasses at $S_d$. Hence the number of proper classes per genus is at most $4\cdot2^{|S_d|}$, with the factor 4 allowing for the dyadic units. This gives a summable majorant after conditioning on $d$, since
\[
 \sum_{d}2^{|S_d|}c(d)/d<\infty .
\]
Indeed all primes outside $S$ dividing such $d$ occur with exponent at least two, and the local polynomial bounds just proved apply. The inequality $\sum_{t\le y}\mu^2(t)3^{\omega(t)}\ll y(1+\log y)^2$ follows already from $3^{\omega(t)}\le d_3(t)$ for squarefree $t$.

Fix $d$ and a tuple of local types at $S_d$. Realize each generator of the product of unit squareclasses at $S_d$ by the residues of odd primes in a reduced arithmetic progression modulo $8\prod_{p\in S_d,\,p>2}p$. CRT supplies the progressions, after excluding finitely many primes. If for each of these finitely many generators some such $q\mid t$ carries the isotropic local type, the proper genus has only one class. In fact the idele supported at $q$ with uniformizer $q$ is trivial in \eqref{eq:spinC} by Lemma~\ref{lem:indlocal}. Dividing it by rational $q$ kills the prescribed unit tuple at $S_d$ and leaves only units at other primes, where those units are in the stabilizer image. Thus all generators vanish. Both actual counts give one class for this tuple.

For a particular generator choose a finite subset $F$ of its prime progression. Among squarefree $t$ coprime to $S_d$, giving three labels to every prime, the asymptotic proportion of labelled tuples using no isotropic label at any $q\in F$ equals
\[
 \prod_{q\in F}\frac{1+2/q}{1+3/q}.
\]
At these primes the local factor $1+3q^{-s}$ is replaced by
$1+2q^{-s}$. Extracting $\zeta(s)^3$ from the two products and
applying the preceding convolution estimate gives the proportion. The proportions tend to zero as $F$ exhausts the progression, by the divergence of the reciprocals of primes in reduced residue classes. A finite union over generators shows that all exceptional labelled tuples at fixed $d$ contribute $o_d(X(\log X)^2)$, even after multiplying by the fixed genus-size bound. Finally truncate the summable majorant in $d$ to pass to the full sum. Thus the excess of either actual count over the genus count is $o(X(\log X)^2)$. Partial summation gives the stated real boundary limit, divergence at $1$, and absolute convergence to its right.
\end{proof}

\subsection{Local equivalence and spinor norms at ramified primes}\label{sec:spincompute}
Local types and transporter spinor norms in \(T(K,L)\) admit
finite congruence descriptions. At a prime \(p^e\Vert m\), list the sublattices in \(L_p/p^eL_p\) of relative index \(p^e\). Take integral bases with determinant \(p^e\) in a common parent basis, for example integral column Hermite forms. At \(p\nmid m\) there is only the parent lattice to consider. Work with an integral nonsingular symmetric Gram matrix of a fixed polarization, including at two.

\begin{lemma}\label{lem:spincompute}
Let \(A,B\) be two such local Gram matrices with the same determinant, and let \(\delta=v_p(\det B)\). Enumerate \(U\in{\rm GL}_n(\mathbb Z/p^{\delta+6}\mathbb Z)\) satisfying \(U^tAU=B\bmod p^{\delta+6}\). Every residue solution has an exact integral isometry congruent to it modulo \(p^3\), and every exact isometry is so obtained up to composition with an intrinsic proper isometry of trivial spinor norm. Hence the full list of signs and spinor classes of ambient transporters between the two embedded lattices is obtained from one such lift per solution.
\end{lemma}
\begin{proof}
Lift \(U\) integrally. Write \(U^tAU=B(I+E)\), so \(E\) has entries in \(p^6\mathbb Z_p\) by the adjugate formula and is self-adjoint for \(B\). The binomial series \(D=(I+E)^{-1/2}\) converges, is self-adjoint, and is the identity modulo \(p^3\). Indeed for odd \(p\) the coefficients are \(p\)-integral (binomial coefficients in a \(p\)-adic integer exponent are integral by continuity). At two their denominators divide \(4^k\) in degree \(k\), by the central binomial coefficient formula. Then \((UD)^t A(UD)=B\). 

If two exact solutions agree modulo \(p^3\), their relative transformation \(g\) is an isometry of the source congruent to the identity modulo \(p^3\). Its binomial square root converges integrally by the same bounds on coefficients, and has determinant \(+1\): its determinant is an orthogonal sign and the root is congruent to the identity modulo \(4\) at two (the linear term has valuation at least two, and higher terms also do). At odd \(p\) it is the identity modulo \(p\). Indeed its adjoint equals its inverse by applying the convergent series to \(g^*=g^{-1}\). Equivalently the square root near the identity is unique, by the binomial identities in the commutative algebra generated by \(g\). Thus \(g\) is a square in the proper group, and its spinor norm is trivial. Passing from basis matrices to the common ambient space only conjugates this intrinsic map. The determinant sign of the ambient transporter equals that of the matrix of the exact isometry because the chosen child bases have equal determinant.
\end{proof}

To compute the squareclass of a lifted transporter use reflection reduction in a rational orthogonal basis of \(V_p\). For each remaining basis vector \(v\), at least one of \(Q(v+gv),Q(v-gv)\) is nonzero, since their sum is \(4Q(v)\). Compute both with increasing precision until a nonzero norm and its squareclass are determined. The corresponding reflection, with the reflection in \(v\) if needed, fixes that coordinate after composition. Continue on its orthogonal complement. The product of the resulting norm classes is the spinor norm, and the search terminates also at \(p=2\), where a unit squareclass is determined modulo eight. This evaluates \(SO(K_p)\) by taking \(A=B\) and keeping positive determinant, and all local transporters by taking each image in its local type.

The idelic quotient reduces to finite-dimensional linear algebra. Include the primes dividing \(2m\det B_L\) in a finite set \(S'\). Rationally clear all valuations of any squareclass idele by a product of primes. Modulo good-place units this represents a class by the finite product of unit squareclasses at \(S'\). Impose the unit tuple of \(-1\) and the images, after valuation clearing, of the stabilizer generators at \(S'\). These generate all relations: rational products used to clear valuations are unique modulo a rational square and a sign, and at primes outside \(S'\) the stabilizer consists of unit squareclasses. This is finite linear algebra over \(\mathbb F_2\). 
For reversal also include the exceptional primes where the chosen rational reflection fails to preserve \(K_p\), then evaluate its products with the improper local lifts. Enlarging \(S'\) in this way gives the same quotient. Products of the local transporter lists and reduction in that finite vector space compute both coefficients for every nondegenerate indefinite parent of rank at least three.


\subsection{Nonmultiplicativity in indefinite signature}
\begin{theorem}\label{th:indeuler}
For every indefinite integral lattice over $\mathbb Q$ of rank $n\ge3$, neither $a_O$ nor $a_{SO}$ is multiplicative.
\end{theorem}
\begin{proof}
Suppose $a_H$ is multiplicative. We first deduce
$a_H(b)=c(b)$ for every $b$. Fix the finite family of local
types at the primes dividing $b$ and enlarge $S$ to include them. Append squarefree $t$ coprime to this enlarged set. By the fixed-$d$ argument of Theorem~\ref{th:indmean}, averaged over labelled $t$ almost every extended genus has just one proper class. Thus the sum of $a_H(bt)$ divided by the sum of $3^{\omega(t)}$ over these $t\le X$ tends to $c(b)$. By \eqref{eq:indsquarefree} multiplicativity would make the ratio equal to $a_H(b)$ for every $X$.

It remains to produce a local tuple containing two represented
ordinary classes. Put $S=\{p:p\mid2\det B_L\}$. Choose a prime $q\notin S$ congruent to $1$ modulo $8$ and modulo every odd prime in $S$. Quadratic reciprocity gives $(\ell/q)=1$ for $\ell\in S$ and $(-1/q)=1$. The determinant of $L_q$ is a unit square, so this unimodular odd-prime lattice admits a basis with Gram matrix the identity. At $q$ replace the basis by $v_1,qv_2,\ldots,q^{n-1}v_n$, and elsewhere keep $L_p$. This gives a child $K\subset L$ of index $b=q^{n(n-1)/2}$. For spinor norms in this paragraph take the rational quadratic form $x\mapsto B(x,x)$; changing this normalization does not change the equivalence relation.

All isometries of $K_q$ have square spinor norm, including improper ones. Its norm coefficients are $1,q^2,\ldots,q^{2(n-1)}$. If $h$ is an integral isometry, the first coordinate of $hv_1$ is $\pm1$ modulo $q$. One of $hv_1\pm v_1$ therefore has unit square norm and its reflection is integral and sends $hv_1$ to $\mp v_1$. Correct the sign with the reflection in $v_1$ if necessary. On the orthogonal complement divide the form by $q^2$ and repeat. This division by a square leaves all reflection norm classes unchanged. The induction proves the assertion.

For a finite idele $x=(x_p)$ consider the quadratic character
\[
 \Xi(x)=\left(\frac{\overline{x_q q^{-v_q(x_q)}}}{q}\right)
          \prod_{\ell\ne q}\left(\frac{\ell}{q}\right)^{v_\ell(x_\ell)} .
\]
It is trivial on squares and on $\mathbb Q^\times$: check rational generators $-1,q,\ell\ne q$. It is also trivial on the spinor norms of the full local stabilizers $O(K_p)$. At $q$ use the preceding induction, at $p\in S$ use $(p/q)=1$, and elsewhere the reflection-generation argument of Lemma~\ref{lem:indlocal} gives even valuation. Since the spinor norm is defined on the full orthogonal group,
$\Xi$ applies to both proper and improper stabilizers.

Choose $g_q\in SO(L_q)$ of nonsquare unit spinor norm, possible by Lemma~\ref{lem:indlocal}, and $g_p=1$ elsewhere. The patched child $K'$ with local lattices $g_pK_p$ still lies in $L$. If a rational isometry $u$ carried $K$ to $K'$, then $(u^{-1}g_p)_p$ would belong to the full local stabilizers of $K$, and $\Xi(\theta(g))$ would be $1$ by the two triviality assertions. By construction it is $-1$. Thus the tuple containing $K,K'$ contributes at least two ordinary classes and hence at least two proper classes. As every other tuple is represented, $a_H(b)>c(b)$, giving the required contradiction.
\end{proof}

\section{The second term and the singularity at one in indefinite signature}\label{sec:spinboundary}
Fix a base index $d$ as in Theorem~\ref{th:indmean}, a local-type
tuple at $d$, and its reference child $K$. Take $C=C(K)$ for proper
equivalence and $C=C(K)/\langle\delta_K\rangle$ for ordinary
equivalence. Let $T\subseteq C$ be the nonempty image of $T(K,L)$;
in general $T$ is a subset rather than a subgroup. Write these
groups additively. For $p\notin S_d$, let $v_p\in C$ be the class of the
idele equal to $p$ at $p$ and to $1$ elsewhere. Characters of $C$ will be denoted by $\chi$.

\begin{lemma}\label{lem:projection}
Give each $p\mid t$ one of the three local labels of Lemma~\ref{lem:indlocal}, and let $D\subset C$ be the subgroup generated by the $v_p$ whose label is isotropic. The number of actual $H$-classes in this extended tuple is $|T\bmod D|$.
\end{lemma}
\begin{proof}
At a nonisotropic label the full stabilizer of the child $U\perp p^2\langle u\rangle$ has spinor norms of even valuation: successively fix a unit diagonal basis of the unimodular $U$ by the unit reflections from Lemma~\ref{lem:indlocal}, then the unary factor also has norm of even valuation. The proper stabilizer contains all units. At an isotropic label its proper stabilizer has full spinor norm, so the new proper genus group is the old one modulo exactly the listed uniformizers. The other local stabilizers have not changed.

At each new prime the transporters of that local type are $SO(L_p)SO(K'_p)$ for the extended child $K'$. Their image is trivial. Thus the represented spinor labels are the quotient image of the old ones. The reversal translation descends to the same quotient: at every new prime choose an improper isometry preserving both the parent and the new child, by reflection in the unit summand perpendicular to the chosen nonisotropic line or to the chosen hyperbolic plane. Such choices can already be used to compute $\delta_K$. This proves the description for either group.
\end{proof}

For a subgroup $W\subset C$ put $\eta=[C:W]^{-1}$. A character of $C$ gives at the good primes a quadratic Dirichlet character $\chi(p)=\chi(v_p)$. These characters factor modulo $Q_d=8\prod_{p\in S_d,\,p>2}p$.
Dividing a uniformizer idele by the corresponding rational prime
leaves unit classes at $S_d$, determined modulo $8$ or $p$. Conversely unit tuples generate $C$ and all unit tuples are attained by prime progressions, so a nontrivial character gives a nonprincipal Dirichlet character. Characters are extended by zero at primes of the modulus. Define
\[
 J_W(s)=\prod_{p\notin S_d}\bigl(1+(2+\mathbf1_{v_p\in W})p^{-s}\bigr),\qquad
 k_W=\lim_{x\downarrow0}x^{2+\eta}J_W(1+x).
\]
The Euler product counts squarefree labelled extensions with
isotropic primes mapping into $W$. Character orthogonality yields
\[
 J_W(s)=\zeta(s)^{2+\eta}
       \prod_{\substack{\chi\in W^\perp\\\chi\ne1}} L(s,\chi)^\eta\,A_W(s).
\]
On the real axis $s>1$ the powers take their positive values. The Euler factors of $A_W$ are the quotients of the corresponding local factors. Outside $S_d$ their linear terms cancel, so the product is absolutely convergent near $1$, with positive value there. At the finitely many discarded primes the residual factor is $(1-p^{-s})^{2+\eta}$. Dirichlet's theorem $L(1,\chi)>0$ for nonprincipal quadratic characters now gives $k_W=A_W(1)\prod_{\chi\ne1}L(1,\chi)^\eta$, with characters restricted as in the display.

\begin{theorem}\label{th:indsecond}
For either orientation group in every indefinite rank $n\ge3$ there is a finite, strictly positive constant
\[
 K_H=\sum_d\frac1d\sum_{\text{base tuples at }d}
  \ \sum_{\substack{\chi\in\widehat C,\ \chi\ne1\\
                      \chi\ \text{nonconstant on }T}} k_{\ker\chi}.
\]
Writing $Z_{\mathrm{gen}}(s)=\sum_m c(m)m^{-s}=\zeta(s)^3V(s)$, one has
\begin{align*}
 Z_H(1+x)-Z_{\mathrm{gen}}(1+x)&\sim K_H x^{-5/2}\quad(x\downarrow0),\\
 \sum_{m\le X}a_H(m)&=\frac{C_L}2X(\log X)^2+
 \frac{K_H}{\Gamma(5/2)}X(\log X)^{3/2}+o_L(X(\log X)^{3/2}).
\end{align*}
Consequently no meromorphic germ at $1$ agrees with $Z_H$ on a right-hand neighbourhood.
\end{theorem}
\begin{proof}
\emph{Separating the first obstruction.}
For fixed $d$ and base tuple expand the function $f(D)=|T\bmod D|$ on the finite set of subgroups as $\sum_{D\subseteq W}\beta(W)$. The coefficients are determined descending from $\beta(C)=1$. For a hyperplane $W=\ker\chi$, $\beta(W)$ is $1$ if $\chi$ takes both signs on $T$, and zero otherwise. Thus the labelled-extension series is $\sum_W\beta(W)J_W(s)$. Subtracting its genus part removes the term $W=C$. Hyperplanes give exponent $5/2$ and the remaining proper subgroups give exponents at most $9/4$. This gives the real-axis asymptotic for each fixed tuple.

\emph{From Euler products to unsmoothed counts.}
For the per-tuple counting limit factor $J_W(s)=\zeta(s)^2 H_W(s)B_W(s)$ where
\[
 H_W(s)=\prod_{\substack{p\notin S_d\\v_p\in W}}(1-p^{-s})^{-1}.
\]
The Dirichlet coefficients of $H_W$ are nonnegative, and $B_W$ has absolutely convergent Dirichlet series on $\Re s>1/2$ by cancellation of the linear terms. At $1$ it is positive. Writing $h(m)$ for the coefficients of $H_W$, the real limit gives $H_W(1+x)\sim [k_W/B_W(1)]x^{-\eta}$. 

Apply Karamata's theorem for a Laplace--Stieltjes transform to the measure with mass $h(m)/m$ at $\log m$ \cite[Thm.~1.7.1]{BGT}. Its cumulative mass up to $\log y$ is asymptotic to $k_W(\log y)^\eta/(B_W(1)\Gamma(1+\eta))$. 

In the convolution with $d_2$, use $\sum_{j\le y}d_2(j)=y\log y+O(y)$ for $y\ge1$. The main sum divided by $y$ is the integral of that cumulative mass from $0$ to $\log y$, while the error divided by $y$ is bounded by the cumulative mass itself. Consequently after convolution with $B_W$ the summatory function of $J_W$ is asymptotic to $k_W y(\log y)^{1+\eta}/\Gamma(2+\eta)$. 
The last convolution follows by dominated convergence, using the resulting bound $O_W(y(1+\log y)^{1+\eta})$ before convolution and absolute convergence at $1$. The resulting estimate holds for each fixed tuple.

\emph{Uniformity over base indices.}
To sum over all base indices $d$, use the pointwise majorant
\[
 0\le |T\bmod D|-1\le |C/D|-1
 =\sum_{\chi\ne1}\mathbf1_{D\subseteq\ker\chi}.
\]
For each such hyperplane the coefficients of $J_{\ker\chi}$ are bounded by those of $\zeta(s)^2H_{\ker\chi}(s)$, by inspecting each nonnegative local factor. On real $1<\sigma\le2$ one has $H_{\ker\chi}(\sigma)^2\le\zeta(\sigma)L(\sigma,\chi)$, locally at split, inert and removed primes. Periodicity and zero mean give $L(\sigma,\chi)\ll1+\log Q_d$ uniformly: sum to $Q_d$ trivially and use partial summation with partial sums bounded by $Q_d$ for the tail. This proves $x^{5/2}J_{\ker\chi}(1+x)\ll\sqrt{1+\log Q_d}$. For summatory coefficients put $\sigma=1+1/\log y$ for $y\ge3$. Then $\sum_{m\le y}h(m)/m\le eH_{\ker\chi}(\sigma)$. The divisor convolution consequently bounds the count for a single hyperplane by $O(y(1+\log y)^{3/2}\sqrt{1+\log Q_d})$, enlarging the constant for small $y$. 
There are at most $4\cdot2^{|S_d|}$ characters. The additional square root of a logarithm is summable against $2^{|S_d|}c(d)/d$: the same Euler majorant converges with $d^{-(1-\epsilon)}$ for $0<\epsilon<1/2$. This justifies both limits by dominated convergence. For the genus term the error $O_L(X(1+\log X))$ follows from the divisor estimate in Theorem~\ref{th:indmean}, since $V$ has absolute convergence to the left of one.

In the construction of Theorem~\ref{th:indeuler}, the two classes in the same local tuple remain distinct even after ordinary orientation reversal. Hence its $T$ has at least two elements in $C$ for either choice of $H$. Some character of this finite binary vector space separates them. The corresponding hyperplane has positive coefficient, proving $K_H>0$. Finally $Z_{\mathrm{gen}}$ is meromorphic at one. If $Z_H$ were meromorphic there, their difference would have an integral order, incompatible with the positive $x^{-5/2}$ asymptotic.
\end{proof}

\subsection{The finite expansions on the real axis}\label{sec:indorders}
Keep the notation of Section~\ref{sec:spinboundary}, and set $\eta_j=2^{-j}$. For each nonempty subset $T$ occurring in a base spinor group $C$ define integers $\beta_T(W)$ by
\[
 |T\bmod D|=\sum_{D\subseteq W\subseteq C}\beta_T(W)
 \qquad (D\subseteq C).
\]
Taking $D=C$ gives $\beta_T(C)=1$; the remaining coefficients follow by downward induction. For $j\ge1$ put
\[
 K_{H,j}=\sum_d d^{-1}\sum_{\text{base tuples at }d}
       \sum_{\substack{W\subseteq C\\\operatorname{codim} W=j}}\beta_T(W)k_W.
\]

\begin{theorem}\label{th:indorders}
For every indefinite parent of rank at least three, every $H=O,SO$ and each fixed integer $J\ge1$, all the indicated coefficient sums converge absolutely and
\[
 Z_H(1+x)=C_L x^{-3}
       +\sum_{j=1}^J K_{H,j} x^{-2-\eta_j}
       +O_{L,J}\bigl(x^{-2-\eta_{J+1}}\bigr)
       \qquad(x\downarrow0).
\]
Here $K_{H,1}=K_H>0$ and $K_{H,2}\le0$. The latter is strictly negative exactly when some base tuple has a codimension-two quotient in which $T$ has two or three images.
\end{theorem}
\begin{proof}
There is a constant $b_j$ depending only on $j$ such that $|\beta_T(W)|\le b_j$ for all $W$ of codimension $j$. Indeed $|T\bmod W|\le2^j$ and all terms in its recursion are indexed by subspaces of the fixed $j$-dimensional quotient. The number of codimension-$j$ subgroups of $C$ is at most $|C|^j$, by counting ordered tuples of dual characters. Subtract from $|T\bmod D|$ its expansion truncated at codimension $J$. The remainder vanishes if $\operatorname{codim}D\le J$. In the other case $D$ lies in at least one subgroup of codimension $J+1$ and the absolute value of the remainder is bounded by
\[
 O_J(|C|^J)\sum_{\operatorname{codim}W=J+1}\mathbf1_{D\subseteq W}.
\]
Since $J\ge1$, the bound also controls $|T\bmod D|$. Summing the indicator for fixed $W$ over labelled good-prime extensions gives $J_W(1+x)$.

Write $Q=Q_d$ and $\eta=\eta_j$ for a fixed codimension $j$. Uniformly for nonprincipal Dirichlet characters $\chi$ modulo $Q$ and $1\le\sigma\le2$ one has
\[
 |L(\sigma,\chi)|\ll1+\log Q,\qquad
 |L'(\sigma,\chi)|\ll(1+\log Q)^2.
\]
Indeed the partial sums of $\chi$ are bounded by $Q$ and its sum over a full period vanishes. Sum up to $Q$ trivially, then use partial summation on the tail with weights $u^{-\sigma}$ and $(\log u)u^{-\sigma}$. The differentiated tail is uniformly convergent down to $\sigma=1$ and is $O(1+\log Q)$. In particular the difference of the endpoint values at $1+x$ and $1$ is $O(x(1+\log Q)^2)$.

Factor $J_W(s)=\zeta(s)^{2+\eta}A_W(s)\prod_{\chi\in W^\perp,\ \chi\ne1}L(s,\chi)^\eta$ as in Section~\ref{sec:spinboundary}. For $1\le\sigma\le2$, $A_W(\sigma)$ is bounded above uniformly, and
\[
 |A_W'(\sigma)|\ll 1+\log Q.
\]
In fact the logarithm and the logarithmic derivative of each eligible local correction are $O(p^{-2\sigma})$ and $O((\log p)p^{-2\sigma})$ respectively, by cancellation of its linear term. These constants can be taken uniform: the label weight is at most three and the sum of the character exponents is less than one. The finitely many smaller primes contribute bounded local factors. At excluded primes the residual factor is $(1-p^{-\sigma})^{2+\eta}\le1$, whose logarithmic derivative is $O((\log p)/(p-1))$. Sum over $p\mid Q$ for the claimed bound.

All the character values $L(\sigma,\chi)$ in this expression are positive on the real interval in question. Using $|u^\eta-v^\eta|\le |u-v|^\eta$ for $u,v\ge0$ and telescoping the product of $2^j-1$ factors gives, for $0<x\le1$ and $j\ge1$,
\begin{align*}
 x^{2+\eta}J_W(1+x)&\ll_j (1+\log Q)^{1-\eta},\\
 |x^{2+\eta}J_W(1+x)-k_W|&\ll_j x^\eta(1+\log Q)^2 .
\end{align*}
For the second estimate, each character factor contributes at most $O(x^\eta(1+\log Q)^{2\eta})$ and the others cost $O_j((1+\log Q)^{1-2\eta})$ in total. The changes from $A_W$ and from $x\zeta(1+x)=1+O(x)$ satisfy the displayed bound as well.

For each extracted $j\ge1$, replacing $d^{-x}x^{2+\eta_j}J_W(1+x)$ by $k_W$ now costs
\[
 O_j\bigl(x^{\eta_j}(1+\log (2dQ))^2\bigr),
\]
using $|d^{-x}-1|\le x\log d$. For every fixed $a,b\ge0$ one has
\[
 \sum_d \frac{c(d)}d\,
    \bigl(4\cdot2^{|S_d|}\bigr)^a(1+\log(2dQ_d))^b<\infty.
\]
Indeed for primes outside the fixed $S$ every nonzero exponent in $d$ is at least two, and the local type count is polynomial in that exponent. The constant multiplier $2^a$ per new prime leaves absolute convergence at every real $\sigma>1/2$ before the logarithmic weight. This also proves absolute convergence of the constants, since $|C|\le 4\cdot 2^{|S_d|}$. The total replacement error for each $j$ in the original series is therefore $O_{L,j}(x^{-2})$. The residual estimate, the first uniform product bound at $j=J+1$, and the same summability with $a=2J+1$ give $O_{L,J}(x^{-2-\eta_{J+1}})$. Finally the $W=C$ terms form the genus series, which equals $C_L x^{-3}+O_L(x^{-2})$ by its absolutely convergent correction to $\zeta(1+x)^3$.

The hyperplane coefficient is the one in Theorem~\ref{th:indsecond}. For codimension two, consider the image of $T$ in $\mathbb F_2^2$. Images of sizes one, two, three and four have respectively zero, two, three and three nonconstant nontrivial characters. Subtract this number and the constant contribution one from the size of the image. The four possible values of $\beta_T(W)$ are $0,-1,-1,0$. All $k_W$ are positive.
\end{proof}

\begin{theorem}\label{th:indorders-terminate}
For the chosen orientation put $C_{SO}(L)=C(L)$, $C_O(L)=C(L)/\langle\delta_L\rangle$, and $r_{H,L}=\dim_{\mathbb F_2} C_H(L)$. Thus $2^{r_{H,L}}$ is the number of $H$-classes in the parent's genus. One has $\beta_T(W)=0$ at every base tuple whenever $\operatorname{codim}W>r_{H,L}+1$. In particular
\[
 Z_H(1+x)=C_Lx^{-3}+\sum_{j=1}^{r_{H,L}+1}K_{H,j}x^{-2-\eta_j}
              +O_L(x^{-2}).
\]
\end{theorem}
\begin{proof}
Let $G_0\subseteq C$ be the image of the spinor norms of $\prod_p SO(L_p)$. The occurring subset satisfies $T+G_0=T$ since postcomposing any local embedding $K_p\to L_p$ by an integral proper parent isometry gives another occurrence. For $H=SO$ take $G=G_0$. The idele formula shows that $C/G$ is a quotient of $C(L)$, since its defining kernel contains every proper local parent stabilizer norm and every rational principal idele.

For $H=O$, choose local improper isometries $i_p$ of $L_p$ and $j_p$ of $K_p$, equal to a fixed rational reflection $r$ at almost all places. For a proper local transporter $g_p$ the tuple $(i_p g_p j_p)_p$ is also admissible, translating its spinor class by $\tau=[\theta((i_p j_p)_p)]$ in $C$. This follows from the multiplicativity of the spinor norm on the full orthogonal group. In squareclasses $\theta(i_pj_p)=\theta(ri_p)\theta(rj_p)$, and $(\theta(rj_p))_p$ is the child reversal, killed in the ordinary quotient. 

Thus $G=G_0+\langle\tau\rangle$ preserves $T$ and $C/G$ is a quotient of the ordinary parent group $C_O(L)$: in addition to all proper parent stabilizer norms its defining kernel kills the idele $(\theta(ri_p))_p$ representing parent reversal. In both cases $P=G^\perp\subseteq\widehat C$ consequently has dimension at most $r_{H,L}$.

For the codimension bound put $A=D^\perp$ and
$R=A\cap P=(D+G)^\perp$: the second space consists of those
characters of $A$ also trivial on the parent-stabilizer image $G$. The image of $T$ in $C/D$ is a union of full cosets of $(D+G)/D$, each of size $|A|/|R|$. The number of these cosets is the size of the image of $T$ in $C/(D+G)$ and depends only on $R$. Thus $|T\bmod D|=|A|h(R)$ for a rational-valued function $h$ on the subspaces of $P$. Explicitly $h(R)=|T\bmod R^\perp|/|R|$. Expand $h(R)=\sum_{Y\subseteq R} b(Y)$ by upward induction in this fixed finite lattice. Since $|A|$ counts the elements of $A$, we obtain the exact identity
\[
 |T\bmod D|=
 \sum_{Y\subseteq P}\ \sum_{\chi\in\widehat C}
 b(Y)\,\mathbf1_{Y+\langle\chi\rangle\subseteq D^\perp}.
\]
The indicator is $\mathbf1_{D\subseteq (Y+\langle\chi\rangle)^\perp}$, and the latter subgroup has codimension at most $r_{H,L}+1$. Uniqueness of the subgroup expansion proves the vanishing. The $O_L(x^{-2})$ estimate follows from the replacement estimates in Theorem~\ref{th:indorders}, now without any unextracted terms.
\end{proof}

\begin{theorem}\label{th:indorders-smooth}
Set $\mathcal B_{H,L}=\{3\}\cup\{2+2^{-j}:1\le j\le r_{H,L}+1\}$. There are real functions $F_{H,a}\in C^\infty([0,1])$ for $a\in\mathcal B_{H,L}$ such that
\[
 Z_H(1+x)=\sum_{a\in\mathcal B_{H,L}} x^{-a} F_{H,a}(x),\qquad 0<x\le1.
\]
Their values at zero are $C_L$ for $a=3$ and $K_{H,j}$ for $a=2+\eta_j$. For each real $b$, the sum of the Taylor terms
$F_{H,a}^{(k)}(0)x^{k-a}/k!$ with $k\ge0$ and $k-a\le b$ differs from $Z_H(1+x)$ by $o_{L,b}(x^b)$ as $x\downarrow0$. The derivative bounds in this real expansion use the ineffective constants in Siegel's theorem.
\end{theorem}
\begin{proof}
The genus term gives $F_{H,3}(x)=(x\zeta(1+x))^3V(1+x)$. For $a=2+\eta_j$ use the real series
\[
 F_{H,a}(x)=\sum_d \frac1d\sum_{\text{base tuples at }d}
  \sum_{\operatorname{codim}W=j}\beta_T(W)\,d^{-x}
             x^{a}J_W(1+x).
\]
By Theorem~\ref{th:indorders-terminate} the identity holds for
$0<x\le1$. We verify termwise differentiation at $x=0$.

For a nonprincipal real character $\chi$ of modulus $Q=Q_d$ and each $\epsilon>0$,
\begin{equation}\label{eq:indorders-siegel}
 \inf_{0\le x\le1} L(1+x,\chi)\gg_\epsilon Q^{-\epsilon}.
\end{equation}
Siegel's theorem gives $L(1,\chi^*)\gg_\lambda q^{-\lambda}$ for the inducing primitive character of conductor $q\mid Q$ and any $\lambda>0$ \cite[Chap.~21]{Davenport}. The imprimitive factors are bounded below by $\prod_{p\mid Q}(1-1/p)\gg_\lambda Q^{-\lambda}$: at each sufficiently large $p$ use $1-1/p\ge p^{-\lambda}$. Thus $L(1,\chi)\ge c_\lambda Q^{-2\lambda}$. With the first-derivative bound in Theorem~\ref{th:indorders}, this gives $L(1+x,\chi)\ge c_\lambda Q^{-2\lambda}/2$ for $0\le x\le c'_\lambda Q^{-2\lambda}(1+\log Q)^{-2}$. For larger $x\le1$, the Euler product gives $L(1+x,\chi)\ge\zeta(2+2x)/\zeta(1+x)\gg x$, also when primes divide the modulus. Taking $\lambda$ sufficiently small in terms of $\epsilon$ gives \eqref{eq:indorders-siegel}.

For $k\ge0$, period summation with weight $(\log u)^k u^{-\sigma}$ gives
$|L^{(k)}(\sigma,\chi)|\ll_k(1+\log Q)^{k+1}$ uniformly on $1\le\sigma\le2$. Indeed the tail beyond $Q$ is bounded by its boundary value times $Q$ and $Q$ times the integral of the absolute derivative of that weight. The resulting tail is $O_k((1+\log Q)^k)$ for $\sigma\ge1$.
The product and chain rules, together with
\eqref{eq:indorders-siegel}, then bound all derivatives through
order $k$ of $L(1+x,\chi)^{\eta_j}$ by
$O_{k,j,\epsilon}(Q^\epsilon)$, taking smaller exponents in
\eqref{eq:indorders-siegel} for the inverse powers.

For $k\ge1$, the eligible-prime terms in the $k$-th derivative of $\log A_W(\sigma)$ are $O_{k,j}((\log p)^k p^{-2\sigma})$ by the same cancelled first-order term. For large primes this follows by termwise differentiation of the Taylor series in $p^{-\sigma}$, uniformly on $1\le\sigma\le2$; small primes contribute boundedly. 

The excluded primes contribute in total $O_{k,j}((1+\log Q)^k)$ by differentiating $\log(1-p^{-\sigma})$ and summing over $p\mid Q$. Together with the upper bound for $A_W$, these estimates bound its ordinary derivatives by fixed powers of $1+\log Q$. The smooth factor $(x\zeta(1+x))^{2+\eta_j}$ and its derivatives are bounded on the real interval. Thus each derivative through order $k$ of $x^aJ_W(1+x)$ is $O_{k,j,\epsilon}(Q^\epsilon)$, uniformly up to $x=0$, by the normalized product formula. In taking products one may shrink $\epsilon$ in each factor, since the number of characters is fixed by $j$.

Differentiating $d^{-x}$ costs powers of $\log d$. The powerful-index summability in Theorem~\ref{th:indorders} and $Q_d\ll_L d$ therefore give uniform absolute convergence of every differentiated series on $[0,1]$, using any sufficiently small final $\epsilon>0$. All the functions therefore extend smoothly to $x=0$ with the
stated derivatives. Taylor's theorem, applied to the finite set
of exponents $a$, gives the expansion.
\end{proof}

\begin{corollary}\label{cor:indorders-oneclass}
For either orientation $K_{H,2}=0$ if and only if the indefinite parent has one $H$-class in its genus. Otherwise $K_{H,2}<0$. In the one-class case $K_{H,j}=0$ for every $j\ge2$, and
\[
 Z_H(1+x)=C_Lx^{-3}+K_Hx^{-5/2}+O_L(x^{-2}).
\]
The one-class case holds in both orientations if the parent contains an integral hyperbolic plane as an orthogonal direct summand.
\end{corollary}
\begin{proof}
If $C_H(L)=0$, the vanishing and the estimate follow from Theorem~\ref{th:indorders-terminate}. Every $H$-class in every child genus is represented. Indeed if the local lattices of a child in this genus are $g_p K_p$ under proper local isometries, patch their overparents $g_p L_p$. The resulting global lattice is $H$-isometric to $L$ by hypothesis, giving the desired embedding, with the chosen equivalence. Thus $T=C$ and $|C/D|$ simply counts all characters trivial on $D$.

Conversely, suppose $C_H(L)\ne0$. Use the prime $q$ and child $K$ of Theorem~\ref{th:indeuler}, leaving all other local lattices equal to $L_p$. The full stabilizer of $K_q$ has square spinor norms, contained in that of $L_q$. There is therefore a surjection $C(K)\to C(L)$. At $q$ reflection in the first basis vector of the chosen unimodular diagonalization preserves both parent and child. At other primes their lattices agree. Taking common improper local symmetries gives $\delta_K\mapsto\delta_L$, hence a surjection in either orientation. 

The image of the whole occurring subset $T$ in $C_H(L)$ is zero. Indeed away from $q$ local transporters are parent stabilizers. At $q$ all units are parent spinor norms. The uniformizer idele at $q$ also dies in $C(L)$: divide by rational $q$, use $q\equiv1\bmod8$ and $q\equiv1\bmod p$ at the odd $p\in S$ to obtain local squares there, and absorb units at all remaining primes. Thus every idele supported at $q$ dies in the parent group. 

By the distinguishing character in Theorem~\ref{th:indeuler}, $T$ has at least two distinct elements even for ordinary equivalence. A nonzero parent character pulls back to a nonzero character constant on $T$, whereas some other character takes both values on $T$. These characters are independent, and the image of $T$ in their rank-two quotient has exactly two elements. The strictness criterion in Theorem~\ref{th:indorders} gives $K_{H,2}<0$.

For a hyperbolic plane with form $xy$, the product of reflections in $(1,1)$ and $(1,u)$ is an integral rotation of spinor norm $u$ for each local unit $u$, including at two. The calculation applies to any common scalar multiple of this plane. Thus all local unit squareclasses occur in the parent stabilizers. Dividing an idele by a rational number to make it a unit everywhere proves that the proper spinor group of the parent is trivial. Eichler's theorem gives the class-number assertion.
\end{proof}

\subsection{The first difference between orientations in indefinite signature}
\begin{theorem}\label{th:indorient}
For every indefinite parent of even rank \(n\ge4\), the constants in Theorem~\ref{th:indsecond} satisfy \(K_{SO}>K_O>0\). In particular
\[
 \sum_{m\le X}\bigl(a_{SO}(m)-a_O(m)\bigr)
 \sim\frac{K_{SO}-K_O}{\Gamma(5/2)}X(\log X)^{3/2}.
\]
In odd rank the two coefficient sequences are identical.
\end{theorem}
\begin{proof}
At each base tuple use the proper group \(C\), its reversal element \(\delta_K\) and its support \(T\). The characters of the ordinary quotient are exactly the proper characters taking value \(+1\) on \(\delta_K\). The nonconstancy condition on \(T\) and the Euler factors
\(J_{\ker\chi}\) are unchanged under this restriction. The difference \(K_{SO}-K_O\) is therefore the absolutely convergent sum of positive terms \(k_{\ker\chi}/d\) over base characters nonconstant on \(T\) and taking value \(-1\) on \(\delta_K\). We construct such a term.

Choose \(q\) as in Theorem~\ref{th:indeuler} and choose a nonsquare unit \(u\in\mathbb Z_q^\times\). As \(n\) is even, the form with Gram \(uI_n\) has square unit determinant, the same squareclass as \(L_q\). Classification of nonsingular symmetric forms over \(\mathbb F_q\) and lifting give a basis of \(L_q\) with this Gram matrix. Starting from the residue-field congruence modulo \(q\), the
matrix square-root correction of Lemma~\ref{lem:spincompute}
gives an integral lift, since \(q\) is odd. In this basis take the child with successive axis scales \(1,q,\ldots,q^{n-1}\) and keep \(L_p\) elsewhere. Write \(K\) for the resulting global child. At \(q\), dividing the form by \(u\) gives the separated diagonal square form of Theorem~\ref{th:indeuler}. Its entire intrinsic orthogonal group has square spinor norm for the divided form. Restoring \(u\) multiplies the spinor norm of a product of \(k\) reflections by \(u^k\), with \((-1)^k\) its determinant. Consequently every proper intrinsic spinor norm at \(q\) is a square, and every improper one is a nonsquare unit modulo squares.

The character \(\Xi\) of Theorem~\ref{th:indeuler} therefore annihilates all proper stabilizer norms of \(K\) and descends to \(C(K)\). In the definition of \(\delta_K\) a rational reflection contributes trivially to \(\Xi\) by rational reciprocity. The local improper stabilizers at \(p\ne q\) also contribute trivially, since \((p/q)=1\) for \(p\in S\), and their spinor norms at the other odd unimodular places have even valuation. The improper stabilizer at \(q\) contributes \(-1\), so \(\Xi(\delta_K)=-1\). An element of \(SO(L_q)\) of nonsquare unit spinor norm takes \(K_q\) to another sublattice in \(L_q\) of the same local type. Extending by the identity at the other places changes the value of \(\Xi\) on the occurring support. Thus this character is nonconstant on the represented support and
takes value $-1$ on $\delta_K$, yielding a positive summand in
$K_{SO}-K_O$. The asymptotic follows by subtracting the two
expansions in Theorem~\ref{th:indsecond}. For odd rank use the central improper isometry \(-I\).
\end{proof}


\begin{corollary}
In even indefinite rank at least four, the set of indices for which \(a_{SO}(m)>a_O(m)\) has a natural density strictly between zero and one. Write \(\mathcal D_*\) for those base indices \(d\) that have at least one tuple with \(\delta_K\ne0\) in its proper group and \(|T|>1\). Its density equals
\[
 \frac1{\zeta(2)}\sum_{d\in\mathcal D_*}\frac1d
                   \prod_{p\in S_d}(1+p^{-1})^{-1}.
\]
\end{corollary}
\begin{proof}
For a fixed tuple the difference counts pairs of distinct represented proper labels identified by reversal after taking the quotient by \(D\). If \(\delta_K=0\) or \(|T|=1\), no extension of this tuple contributes. Otherwise take two labels with difference \(w\ne0\). For \(w=\delta_K\), choosing the nonisotropic labels at every new prime already gives a pair. If \(w\ne\delta_K\), then \(w+\delta_K\ne0\) and its span does not contain \(\delta_K\). There is a prime progression of positive density on which \(v_p=w+\delta_K\), by the unit generation and congruence description preceding Theorem~\ref{th:indsecond}. If \(t\) contains a prime in this progression, give that prime the isotropic label and all others nonisotropic labels. The two distinct quotient labels differ by the nonzero image of \(\delta_K\), giving the desired pair.

Among squarefree \(t\) coprime to \(S_d\), the proportion divisible by no prime in this progression is zero. Indeed for a finite subset of its primes the proportion avoiding it tends to the product of \((1+1/p)^{-1}\) over that subset, by the elementary squarefree density formula. These products tend to zero by divergence of the reciprocal primes in the progression. Thus the conditional proportion with an orientation gap tends to one for each \(d\in\mathcal D_*\) and is zero otherwise. The density of the component indexed by \(d\) is the displayed weight, and the upper bound \(X/d\) allows dominated convergence since \(\sum_{d\in\mathcal D}1/d<\infty\). The constructed nonsquare diagonal tuple belongs to \(\mathcal D_*\), whereas \(d=1\) does not. Both it and the complement have positive weight.
\end{proof}

\subsection{Density of indices with multiple represented classes per local type}
Retain \(S=\{p:p\mid 2\det B_L\}\) in the indefinite setting \(n\ge3\).
For every \(m\), write \(m=dt\), where the entire \(S\)-part and every other prime power of exponent at least two belong to \(d\). The remaining \(t\) is squarefree and coprime to \(S_d\). Let \(\mathcal D\) be the set of admissible \(d\), including one. Write \(\xi_H(m)=1\) if \(a_H(m)>c(m)\), and zero otherwise.

\begin{theorem}\label{th:densityind}
One has \(\xi_H(m)=\xi_H(d)\). The set of indices with more represented \(H\)-classes than local type tuples has a natural density strictly between zero and one. More precisely,
\[
 \sum_{m\le X}\xi_H(m)=\Delta_H X+O_{L,\epsilon}(X^{1/2+\epsilon}),\qquad
 \Delta_H=\frac{1}{\zeta(2)}
 \sum_{\substack{d\in\mathcal D\\a_H(d)>c(d)}}\frac1d
                 \prod_{p\in S_d}(1+p^{-1})^{-1}.
\]
\end{theorem}
\begin{proof}
If each base tuple at \(d\) represents a single spinor label in its relevant quotient, any further quotient of that singleton has one element. If a tuple has more than one label, choose either nonisotropic extension at every prime dividing \(t\). For all \(2^{\omega(t)}\) choices the subgroup \(D\) in Lemma~\ref{lem:projection} is zero; the labels of this tuple remain distinct. Every other tuple contributes at least one class. Hence $\xi_H(m)=\xi_H(d)$.

For \(\Re s>1\) sum the squarefree extensions for each admissible \(d\). Their Euler product gives
\[
 \sum_{m\ge1}\xi_H(m)m^{-s}=\zeta(s) U_H(s),\qquad
 U_H(s)=\frac1{\zeta(2s)}\sum_{d\in\mathcal D}\xi_H(d)d^{-s}
                 \prod_{p\in S_d}(1+p^{-s})^{-1}.
\]
The latter is an absolutely convergent Dirichlet series on \(\Re s>1/2\). Indeed the finite inverse factors have absolute coefficient norm at \(\sigma>1/2\) bounded by \(O_{S,\sigma,\eta}(d^\eta)\) for any \(\eta>0\), and \(\sum_{d\in\mathcal D}d^{-u}\) converges for \(u>1/2\) by its elementary prime product. The reciprocal of \(\zeta(2s)\) is also absolutely convergent there. In the convolution with the constant one sequence, summing \(\lfloor X/k\rfloor=X/k+O(1)\) gives the stated error and constant. The positive spinor twist in Theorem~\ref{th:indeuler} supplies at least one admissible \(d\) of positive weight with \(\xi_H(d)=1\); its prime exponent is \(n(n-1)/2\ge3\). Taking \(d=1\) supplies strictly positive weight to the complement, whose density is given by the same proof with \(1-\xi_H\).
\end{proof}
\begin{remark}
Both indicator Dirichlet series extend meromorphically to
$\Re s>1/2$, with positive simple-pole residues at $s=1$.
\end{remark}


\section{Indefinite binary parents}\label{sec:realbinary}
The rationally split indefinite binary case was treated by Kim,
Lee and Lee \cite{Hyperbolic}; its proper integral stabilizers are
finite. For a nonsplit space, the proper rotation group modulo
scalar signs is infinite cyclic, and the order of its generating
quadratic unit already enters the prime-index count.

\begin{proposition}\label{prop:splitbinary}
For a rationally split indefinite binary parent the restoration formula, the meromorphic continuation and error in Theorem~\ref{th:analytic}, and the residue comparison in Theorem~\ref{th:orientation} all hold with \(n=2\). The proper transporter sum has only its scalar terms.
\end{proposition}
\begin{proof}
The two rational isotropic lines span the space. Any proper isometry preserves both lines, since in an isotropic basis the isometries interchanging the lines are antidiagonal of determinant \(-1\). An integral automorphism acts by signs on the primitive generators of the intersections of those lines with its lattice. Its determinant being one, the signs coincide. Thus every proper intrinsic stabilizer is \(\{\pm I\}\) and the ordinary stabilizer has order at most four. The finite stabilizers allow the restoration and HNF estimates of Theorems~\ref{th:core} and~\ref{th:analytic} to apply in rank two. The orientation comparison uses the same scalar masses and a rational orthogonal basis. Integral equivalence is decided by the finite procedure below.
\end{proof}

Both equivalence relations admit finite tests at every index.
Start from column HNFs and their Gram matrices. To test \(U^t A U=B\), make integral proper basis changes so \(a=A_{11}\) and \(b=B_{11}\) are nonzero (one of three primitive directions suffices). Write \(A=\left(\begin{smallmatrix}a&d\\d&c\end{smallmatrix}\right)\), \(D=d^2-ac>0\). For the first column \((u,v)^t\), put \(t=au+dv\); then \(t^2-Dv^2=ab\ne0\). For square \(D\) factor \(ab\) into the two signed integer factors \(t\pm\sqrt D\,v\). For nonsquare \(D\), take any positive Pell unit \(\lambda=T+S\sqrt D>1\) of norm one, available by the classical Pell theorem. Multiplication by this unit acts on \((u,v)^t\) by the proper integral isometry
\[
 \begin{pmatrix}T-dS&-cS\\ aS&T+dS\end{pmatrix}.
\]
Modulo its powers it suffices to enumerate the solutions with
\(\sqrt{|ab|}\le |t+v\sqrt D|<\lambda\sqrt{|ab|}\). The conjugate is \(ab/(t+v\sqrt D)\), giving explicit bounds for both integers \(t,v\). Discard pairs with \(a\nmid t-dv\). For each remaining first column and each desired sign, the two linear equations for the second column \(w\) are \(x^tAw=B_{12}\) and \(\det(x,w)=\pm1\), with \(x=(u,v)^t\). Their coefficient determinant is \(b\ne0\). Test the unique solution for integrality and for \(w^tAw=B_{22}\). The resulting finite tests decide ordinary and proper equivalence.

\begin{theorem}\label{th:inertreal}
Suppose the indefinite binary space is rationally nonsplit. For every index \(m\) supported on odd primes \(p\nmid\det B_L\) with \((-\det B_L/p)=-1\), the following formula computes both coefficients from a single generator of the rotation group modulo \(\{\pm I\}\). For each \(p^e\Vert m\), independently choose \(k_p=e,e-2,\ldots\) down to \(0\) or \(1\). Set \(v_p(0)=o_p(0)=1\). For \(k>0\) put
\[
 v_p(k)=(p+1)p^{k-1},\qquad
 o_p(k)=\operatorname{ord}\bigl(\varepsilon\text{ in }
                     {\rm PGL}(L_p/p^k L_p)\bigr).
\]
For the tuple $\boldsymbol{k}=(k_p)_{p\mid m}$ define
\(\mathfrak a(\boldsymbol{k})=\prod_p v_p(k_p)/\operatorname{lcm}_p o_p(k_p)\),
where the least common multiple of an empty set is one, and let
\(r=\#\{p:k_p>0\}\).
The proper coefficient is the sum of $\mathfrak a(\boldsymbol{k})$ over these tuples. If the ordinary parent has an improper automorphism, replace this weight by
$[\mathfrak a(\boldsymbol{k})+2^r]/2$; otherwise the ordinary coefficient equals the proper coefficient at these indices.
\end{theorem}
\begin{proof}
Over the reals the proper orthogonal group in split coordinates is diagonal, with weights \(z,z^{-1}\). Its integral stabilizer modulo scalar signs embeds discretely in \(\mathbb R\) by \(\log|z|\), since bounded intervals give bounded integral matrices. The Pell construction supplies a nontrivial element, so the image is infinite cyclic; choose a generator \(\varepsilon\). A generator is obtained by enumerating integral rotations whose
logarithms lie between zero and that of a known Pell element.

At an indicated prime the reduced form is anisotropic and unimodular. Thus the valuation of the norm of any vector is twice the minimum valuation of its coordinates, up to the fixed unit normalization of the polarization. In particular \(L_p=\{x:Q(x)\in\mathbb Z_p\}\) in this normalization, and every rational local isometry preserves \(L_p\). At primes away from \(m\) the children coincide with \(L_p\). Every isometry between children of such indices therefore preserves the global parent. By Smith normal form, removing the common scalar \(p^{(e-k_p)/2}\) leaves the kernel of a primitive covector modulo \(p^{k_p}\), or all of \(L_p\) for \(k_p=0\). There are \(v_p(k_p)\) choices.

If a rotation fixes one primitive projective covector modulo \(p^k\), its corresponding normal has unit norm. The eigenvalue has square one, hence equals \(+1\) or \(-1\) since \(p\) is odd. The complementary line is nondegenerate and its eigenvalue is the same sign by determinant one. Thus the rotation is projectively scalar on the entire module. The simultaneous rotation orbit has length exactly the stated least common multiple. If an improper parent automorphism exists, it and every element of its rotation coset have square one and determinant minus one in the rational split eigenbasis. Their local projectors \((1\pm h)/2\) have rank one, hence precisely two primitive projective lines are fixed at every active prime power. Taking the average over one period of the rotations and one period of their improper coset gives $[\mathfrak a(\boldsymbol{k})+2^r]/2$, also when the projective action is not faithful. Burnside averaging gives the stated ordinary and proper coefficients.
\end{proof}
At a good inert prime the proper coefficient is \((p+1)/o_p(1)\).
For indices with several prime divisors, the local periods enter
through their least common multiple rather than separate factors.


\subsection{Euler products for indefinite binary spaces}
\begin{theorem}\label{th:binaryeuler}
For every indefinite binary parent the ordinary coefficient is nonmultiplicative.
For a rationally nonsplit parent the proper coefficient is nonmultiplicative as well.
For a rationally split parent the proper coefficient is multiplicative precisely when
the positive integer \(B\) of \cite[Thm.~1.1]{Hyperbolic} divides \(24\).
In each multiplicative case
\[
 Z_{SO}(s)=\zeta(s-1)\sum_{h\mid B}h^{-s}.
\]
Here multiplicativity is equivalent to the normalized Euler product of the prime-power subseries.
\end{theorem}
\begin{proof}
First assume the space is nonsplit, and let \(K=\mathbb Q(\sqrt{-\det B_L})\).
Let \(z,z^{-1}\) be the eigenvalues of the generator of the integral rotations
modulo scalar signs in Theorem~\ref{th:inertreal}. They are quadratic integral units
in \(K\), by the diagonal action on the two isotropic lines over \(K\).
Galois conjugation interchanges them. Put \(\alpha=z^2\).
At every good inert odd \(p\) outside the finitely many primes dividing
the discriminant of the characteristic polynomial of the generator,
the two reduced eigenvalues are distinct in \(\mathbb F_{p^2}\).
Consequently the projective order \(o_p=o_p(1)\) is the multiplicative order
of \(\alpha\) in that residue field. Choose such a \(p\) with \(o_p>1\);
infinitely many inert primes are available, whereas \(z^2-1\) is a nonzero
algebraic integer.

Choose a prime \(\ell\mid o_p\), let \(a=v_\ell(p^2-1)\), and put \(N=\ell^a\).
We obtain infinitely many further good inert primes \(q\) with
\(\ell\mid o_q\) by Chebotarev.
In the cyclic group \(\mathbb F_{p^2}^{\times}\) an element is an \(N\)th power
if and only if its order is prime to \(\ell\). Thus \(X^N-\alpha\) has no
root over \(\mathbb F_{p^2}\). Consider the finite extension
\[
 F=K(\zeta_{\ell N},\sqrt[N]{\alpha}).
\]
It is Galois over \(\mathbb Q\): it contains all the roots of \(X^N-\alpha\)
and the inverses of those roots give all the roots of \(X^N-\alpha^{-1}\).
The prime \(p\) is unramified in \(F\): it is unramified in
\(K\), is prime to \(\ell N\), and \(N\) and \(\alpha\)
are local units. The roots consequently generate an unramified
algebra by separable reduction.
Choose a prime of \(F\) above \(p\) with arithmetic Frobenius \(\sigma\).
Its restriction to \(K\) is nontrivial. Since the roots have distinct
reductions, \(\sigma^2\) has no fixed point on the roots of \(X^N-\alpha\).

The Chebotarev density theorem \cite{StevenhagenLenstra} gives
infinitely many rational primes \(q\) for which one can choose a
prime of \(F\) with Frobenius equal to \(\sigma\).
Exclude \(p\) and the finitely many exceptional primes.
The action on \(K\) makes \(q\) inert, and its action on \(\zeta_{\ell N}\)
gives \(q\equiv p\bmod\ell N\), hence \(v_\ell(q^2-1)=a\).
The same permutation \(\sigma^2\) on the roots shows that \(X^N-\alpha\)
has no root over \(\mathbb F_{q^2}\). Cyclicity again yields \(\ell\mid o_q\).

Now \(g=\gcd(o_p,o_q)\ge2\). Write \(A=(p+1)/o_p\) and \(C=(q+1)/o_q\).
The proper coefficients at \(p,q,pq\) are \(A,C,gAC\), by Theorem~\ref{th:inertreal}.
If the parent has no improper integral isometry these are also the ordinary coefficients.
Otherwise the ordinary coefficients are \((A+2)/2,(C+2)/2,(gAC+4)/2\).
Both \(A,C\) are at least two: at an odd inert prime \(z\) reduces to the norm-one
group of order \(p+1\), and \(z^2\) has order dividing \((p+1)/2\).
Four times the ordinary discrepancy at \(pq\) is
\[
 (2g-2)AC+(A-2)(C-2)>0 .
\]
Neither coefficient sequence is therefore multiplicative in the nonsplit case.

In the split plane every integral stabilizer is finite by
Proposition~\ref{prop:splitbinary}. The prime-extension limit of
Lemma~\ref{lem:pr} and its proof use the same local restoration and
finite stabilizers, so they apply here as well.
If the ordinary coefficient were multiplicative, the two large-prime
limits in the proof of Theorem~\ref{th:nonmult} would force
\(|\Gamma_O(L)|/2=1\) and then \(|\Gamma_O(M)|/2=1\) for every represented \(M\).
But a rational orthogonal basis gives a diagonal sublattice inside \(L\),
and changing one coordinate sign is a nonscalar automorphism of that lattice.
The resulting nonscalar automorphism contradicts ordinary
multiplicativity.

For the split proper coefficient use the normalization of
\cite[Thm.~1.1]{Hyperbolic}: embed a scalar copy of \(L\) into the plane
\(H=\mathbb Z^2\) with quadratic form \(xy\), and define
\(B\mathbb Z=[H:L](\mathfrak nL)^{-1}\), where \(\mathfrak nL\) is the ideal
generated by the integral values of \(xy\) on \(L\).
For \(b\mid B\) put
\[
 C_b(m)=\#\{d^2\bmod b:d\mid m,\ (d,b)=1\}.
\]
Their formula is
\[
 Z_{SO}(s)=\frac{\zeta(s-1)}{\zeta(s)}
     \sum_{b\mid B}(B/b)^{-s}\sum_{m\ge1}C_b(m)m^{-s}.
\]
If the unit group modulo \(B\) has exponent dividing two, all units
modulo every divisor \(b\) have square one, by surjectivity of reduction
on unit groups. Then \(C_b(m)=1\) for every \(b,m\); the asserted formula
follows, and its coefficients are multiplicative.
Otherwise take a unit \(v\bmod B\) with \(v^2\ne1\), and choose
distinct primes \(p,q\equiv v\bmod B\). Among indices coprime to \(B\)
only the undilated term \(b=B\) contributes. Thus on \(1,p,q,pq\)
the coefficients are the Dirichlet convolution of Euler's totient with \(C_B\).
Here \(C_B(p)=C_B(q)=2\), and \(C_B(pq)\le3\) since the only residues
are \(1,v^2,v^4\). Multiplicativity of the totient gives
\(a_{SO}(pq)-a_{SO}(p)a_{SO}(q)=C_B(pq)-4<0\).
The unit group modulo \(B\) has exponent at most two exactly when \(B\mid24\).
For an odd prime power \(r^e\mid B\), reduction modulo \(r\) forces \(r=3\);
the unit \(1+3\) rules out \(e\ge2\). At two every odd square is
one modulo eight and \(3^2\ne1\bmod16\). The Chinese remainder theorem combines these local restrictions.
\end{proof}

\subsection{Growth bounds for nonsplit indefinite binary spaces}
\begin{theorem}\label{th:realgrowth}
For a fixed nonsplit indefinite binary parent and either orientation, for every fixed \(\epsilon>0\),
\[
 X^{3/2-\epsilon}\ll_{L,\epsilon}\sum_{m\le X}a_H(m)\ll_L X^2/\log X\qquad(X\ge3).
\]
The abscissa lies in \([3/2,2]\). Along real \(s=2+x\) the series is \(O_L(1+\log(1/x))\) as \(x\downarrow0\), so it cannot have a pole at two.
\end{theorem}
\begin{proof}
\emph{Upper bound.}\quad Realize the space as a line over the real quadratic field \(F\) generated by the parent rotation \(\varepsilon\) of Theorem~\ref{th:inertreal}, acting as multiplication by a norm-one unit. Its eigenvalues are distinct and of absolute values \(\lambda,\lambda^{-1}\), where \(\lambda>1\). Counting orbits of this rotation on raw subgroups of index at most \(X\) gives an upper bound for the number of intrinsic classes. Those of length at least \(J\) contribute \(O(X^2/J)\) by the binary HNF divisor sum. 
Every shorter orbit consists of points fixed by some \(\varepsilon^j\), \(1\le j\le J\). For a lattice \(M\) fixed by that power put \(I=O_F M\subset O_F L\). There are \(O_L(X\log(2X))\) possible \(I\): their index in \(O_F L\) is at most \([O_F L:L]X\) and invertible ideal arithmetic reduces this count to integral ideals, bounded coefficientwise by the divisor function by prime splitting. Since \(\varepsilon^j\) is not rational, \(R_j=\mathbb Z[\varepsilon^j]\) is a quadratic order and \(f_j=[O_F:R_j]\ll_L \lambda^{O(j)}\) by its discriminant \((\varepsilon^j-\varepsilon^{-j})^2\). We have \(f_jI\subset M\subset I\). The number of intervening additive subgroups is bounded by a fixed power of \(1+f_j\), for instance by counting all HNFs of indices at most \(f_j^2\). 
Thus the total short-orbit contribution is \(O_L(e^{c_LJ}X\log(2X))\). Take \(J\) a sufficiently small positive multiple of \(\log X\).

\medskip
\noindent\emph{Lower bound.}\quad Let \(\epsilon>0\) and take \(0<\delta<\min(\epsilon/8,1/16)\). Choose a fixed modulus \(M_0\) divisible by eight and the discriminant of \(F\), so primes \(p\equiv1\bmod M_0\) split in \(F\). Enlarge \(M_0\) to include all bad primes. For large \(y\) consider primes \(q\) in
\([y^{1/2-2\delta},y^{1/2-\delta}]\). There are \(\gg_{L,\delta}y/\log y\) distinct split primes \(p\in(y/2,y]\) for which some such \(q\) divides \(p-1\). In fact sum the primes in residue class \(1\bmod M_0q\) over these \(q\). By the Bombieri--Vinogradov theorem \cite{IK}, applied at both \(y\) and \(y/2\) and using partial summation, the total error from the expected prime counts is \(O_{L,A}(y/(\log y)^A)\) for any fixed \(A\). All moduli are within its range for large \(y\), since \(M_0\) is fixed. The main term is \((\operatorname{Li}(y)-\operatorname{Li}(y/2))\sum_q1/\varphi(M_0q)\), of order \(y/\log y\) by Mertens' theorem over the indicated interval of prime \(q\). An individual \(p\) is counted only \(O_\delta(1)\) times since the product of its distinct divisors \(q\) is at most \(p-1\). This proves the assertion.

Put \(z=\lceil y^{1/2+2\delta}\rceil\). For each of these split primes, \(p-1=q h\) with \(q,h\le z\), so \(p-1\) divides \(\operatorname{lcm}(1,\ldots,z)^2\). Choose products \(m\) of exactly \(k=\lfloor y^{1/2+3\delta}\rfloor\) distinct such primes, and put \(X_y=y^k\). At each \(p\mid m\) there are \(p-1\) primitive normals of nonzero norm, by the split binary finite-field formula. Every isometry between the resulting regular kernels stabilizes the restored parent. The projective parent rotation group is cyclic. Its local projective order divides \(p-1\): a proper split binary rotation preserves both isotropic lines and is diagonal with entries \(t,t^{-1}\) over \(\mathbb F_p\). The simultaneous orbit length is therefore at most \(\operatorname{lcm}(1,\ldots,z)^2\). Consequently either orientation counts at least
\[
 \frac{(y/2-1)^k}{2\operatorname{lcm}(1,\ldots,z)^2}
   =X_y^{1-o(1)}
\]
classes at each of these indices. The ordinary count is always at least half the proper one. Here even \(\log\operatorname{lcm}(1,\ldots,z)\le z\log z=o(k\log y)\) suffices. Unique factorization distinguishes the indices. Their number is at least
\(\binom{\lfloor c_{L,\delta}y/\log y\rfloor}{k}
 \ge X_y^{1/2-3\delta-o(1)}\).
They all lie below \(X_y\), and \(\log X_{y+1}/\log X_y\to1\) for integer \(y\). Monotonicity extends the lower bound to all real \(X\) with the stated loss \(\epsilon\). Partial summation proves the assertions about the abscissa and real boundary growth. Positivity and Landau's theorem show that, if the abscissa is two, this real boundary point is singular despite the absence of a pole there.
\end{proof}



\begin{corollary}\label{cor:real-large-moduli}
For every fixed nonsplit indefinite binary parent, either orientation and every $\epsilon>0$,
\[
 \sum_{m\le X}a_H(m)\gg_{L,\epsilon} X^{13/8-\epsilon}.
\]
In particular the abscissa of convergence is at least $13/8$.
\end{corollary}
\begin{proof}
Let $\vartheta=5/8$ and let $M_0$ be the fixed splitting modulus of Theorem~\ref{th:realgrowth}. Pascadi's estimate for triply-well-factorable weights \cite[Def.~1.1 and Thm.~1.3(i)]{Pascadi} gives, for any $A,\kappa>0$, fixed residue $a=1$ and $Q\le x^{\vartheta-\kappa}$,
\[
 \sum_{d\le Q}\lambda_d\left(\pi(x;d,1)-\frac{\pi(x)}{\varphi(d)}\right)
       \ll_{A,\kappa}\frac{x}{(\log x)^A}
\]
when $\lambda$ is triply-well-factorable of level $Q$. This means that for each $Q=Q_1Q_2Q_3$ with $Q_i\ge1$, the same sequence is the Dirichlet convolution of three $1$-bounded sequences supported on integers $\le Q_i$.

Fix $0<h<\min(\epsilon/20,1/100)$. For large $y$ put $u=y^h$, $Q=y^{\vartheta-h}$ and $N=\lfloor(\vartheta-5h)/h\rfloor$. Divide $(u/2,u]$ into $N$ disjoint equal-length intervals. Let $\lambda$ be the Dirichlet convolution of the point mass at $M_0$ and the $N$ prime indicators of these intervals. Every convolution of a subset of the factors is $1$-bounded by unique factorization, since the prime intervals are disjoint and the factor $M_0$, when present, is fixed. For every $Q=Q_1Q_2Q_3$ the sum of the integer capacities $\lfloor\log Q_i/\log u\rfloor$ is at least $\vartheta/h-4\ge N+1$. Allocate the $N+1$ factors among the three groups within those capacities. Every factor has support at most $u$, including $M_0$ for large $y$, and an empty group is the point mass at one. This proves the required triple factorability for every factorization of $Q$.

Apply the estimate at $x=y$ and $x=y/2$ with $\kappa=h/2$;
the chosen exponent of $Q$ permits both applications for large $y$. The sum of the main coefficients satisfies
\[
 \sum_d\frac{\lambda_d}{\varphi(d)}
 \ \ge\ \sum_d\frac{\lambda_d}{d}
 \ =\ \frac1{M_0}\prod_{i=1}^N\sum_{\ell\ \mathrm{prime\ in\ interval}\ i}\ell^{-1}
 \ \gg_{L,h}(\log y)^{-N},
\]
by the prime number theorem on each fixed-proportion interval. Taking $A>N+1$ supplies $\gg_{L,h}y/(\log y)^{N+1}$ incidences of primes $y/2<p\le y$ with $d\mid p-1$ and $\lambda_d=1$. Each integer $p-1$ has $O_h(1)$ distinct prime divisors exceeding $u/2$, so an individual $p$ accounts for $O_h(1)$ incidences. There are therefore $\gg_{L,h}y/(\log y)^{N+1}$ distinct such primes $p$, all $1\bmod M_0$.

For each of these primes, $p-1=M_0\ell_1\cdots\ell_N b$ with $b\ll_{L,h}y^{1-Nh}$ and $Nh\ge\vartheta-6h$. All $N+2$ factors are bounded by $z=\lceil y^{1-\vartheta+7h}\rceil$ for large $y$. Hence $p-1$ divides $\operatorname{lcm}(1,\ldots,z)^{N+2}$. Set $k=\lfloor y^{1-\vartheta+8h}\rfloor$, $X_y=y^k$. The cyclic-kernel argument of Theorem~\ref{th:realgrowth} gives at least $(y/2-1)^k/(2\operatorname{lcm}(1,\ldots,z)^{N+2})=X_y^{1-o(1)}$ classes at each product of $k$ distinct such primes. Their number of distinct products is at least $X_y^{\vartheta-8h-o(1)}$, since $(\log y)^{N+1}=y^{o(1)}$ for fixed $h$. Monotonicity and $\log X_{y+1}/\log X_y\to1$ give the bound for every sufficiently large real $X$. Positivity and partial summation give the abscissa bound.
\end{proof}

\begin{remark}
The same argument with Maynard's estimate
\cite[Thm.~1.1]{Maynard} and $\vartheta=3/5$ gives the exponent
$8/5-\epsilon$.
\end{remark}
\section{A uniform parity criterion in every signature}\label{sec:allparity}
\begin{theorem}\label{th:allparity}
Let $L$ be a nondegenerate integral quadratic lattice of rank $n\ge1$. Its ordinary and proper coefficient sequences coincide at every index if and only if $n$ is odd. If $n$ is even, there exist an index \(b\) and a finite set of primes \(S'\) depending on $L$ such that
\[
 a_{SO}(bt)>a_O(bt)\qquad\bigl((t,\prod_{p\in S'}p)=1\bigr).
\]
For a definite parent of even rank \(n\ge4\) the difference on this set is \(\gg_L t^{n-1}\); for \(n=2\) it is \(\gg_L\varphi(t)\).
\end{theorem}
\begin{proof}
Odd rank follows by multiplication by \(-I\). For even rank, choose \(q\equiv1\bmod 8\) and \(q\equiv1\bmod p\) at all odd primes dividing the determinant of an integral polarization. The local unimodular Gram has square determinant at \(q\), hence in even rank is isometric to \(uI_n\) for a nonsquare unit \(u\). Take \(K_q\) with scales \(1,q,\ldots,q^{n-1}\) in this basis and keep the parent at other primes. Any improper intrinsic isometry would have nonsquare unit spinor norm at \(q\), even spinor valuation at good primes other than \(q\), and arbitrary norms only at primes where \(\Xi\) is trivial. Applying the explicitly defined character \(\Xi\) of Theorem~\ref{th:indeuler} would then give \(-1\) on a rational spinor norm. Its triviality on rational numbers rules this out. The same spinor-norm obstruction applies in rank two.

Fix any rational reflection \(r\) and scale \(K\) by a single positive integer so that both \(K\) and \(rK\) lie in \(L\). They have the same index \(b\). No proper isometry can take \(K\) to \(rK\), since composing with \(r^{-1}\) would give an improper automorphism of \(K\). Let \(S'\) contain all primes of \(2b\det B_L\). For every \(t\) allowed in the statement, insert at each \(p^e\Vert t\) any nonisotropic primitive kernel in \(K_p=L_p\) and take \(M\) to be the global child. Such kernels exist in every rank at least two by the formula for \(P_n(p)\). Both \(M\) and \(rM\) are contained in \(L\). Restoring their regular parts at all primes dividing \(t\) recovers respectively \(K\) and \(rK\); the uniqueness in Lemma~\ref{lem:local} is purely local and holds in any signature. Every proper identification of these two children would extend to a proper identification of those two different orientations of \(K\). Thus their ordinary class contributes one to the difference.

In definite signature, distinct \(\Gamma_O(K)\)-orbits of such kernels give distinct ordinary classes by the same restoration argument. Their number is at least
\[
 |\Gamma_O(K)|^{-1}t^{n-1}\prod_{p\mid t}
   \left(1-\Bigl(\frac{(-1)^{n/2}\det B_L}{p}\Bigr)p^{-n/2}\right).
\]
For \(n\ge4\) the product is bounded below by a positive constant, and for \(n=2\) it is bounded below by \(\varphi(t)/t\). Negative definite forms may be multiplied by minus one. These bounds give the stated growth estimates.
\end{proof}


\begin{thebibliography}{99}
\bibitem{AGM} D. Allcock, I. Gal and A. Mark, The Conway--Sloane calculus for $2$-adic lattices, \emph{Enseign. Math.} \textbf{66} (2020), 5--31. \url{https://ems.press/journals/lem/articles/17165}.
\bibitem{Baake} M. Baake, Solution of the coincidence problem in dimensions $d\le4$, in \emph{The Mathematics of Long-Range Aperiodic Order} (R. V. Moody, ed.), Kluwer, 1997, 9--44. \url{https://doi.org/10.1007/978-94-015-8784-6_2}.
\bibitem{BaakeZeiner} M. Baake and P. Zeiner, Geometric enumeration problems for lattices and embedded $\mathbb Z$-modules, in \emph{Aperiodic Order} (M. Baake and U. Grimm, eds.), Cambridge University Press, 2017, 73--172. \url{https://doi.org/10.1017/9781139033862.005}.
\bibitem{BGT} N. H. Bingham, C. M. Goldie and J. L. Teugels, \emph{Regular Variation}, Cambridge University Press, 1987. \url{https://doi.org/10.1017/CBO9780511721434}.
\bibitem{ChanOh} W. K. Chan and B.-K. Oh, Can we recover an integral quadratic form by representing all its subforms?, \emph{Adv. Math.} \textbf{433} (2023), 109317. \url{https://doi.org/10.1016/j.aim.2023.109317}.
\bibitem{Chenevier} G. Chenevier, Statistics for Kneser $p$-neighbors, \emph{Bull. Soc. Math. France} \textbf{150} (2022), no.~3, 473--516. \url{https://doi.org/10.24033/bsmf.2852}.
\bibitem{Davenport} H. Davenport, \emph{Multiplicative Number Theory}, 3rd ed., Springer, 2000.
\bibitem{IK} H. Iwaniec and E. Kowalski, \emph{Analytic Number Theory}, AMS Colloquium Publications, vol.~53, American Mathematical Society, 2004. \url{https://doi.org/10.1090/coll/053}.
\bibitem{KLL} D. Kim, S. H. Lee and S. Lee, Zeta functions enumerating subforms of quadratic forms, \emph{Math. Ann.} \textbf{393} (2025), 1459--1514. \url{https://doi.org/10.1007/s00208-025-03273-9}.
\bibitem{KLLpre} D. Kim, S. H. Lee and S. Lee, Zeta functions enumerating subforms of quadratic forms, arXiv:2409.05625v1 (2024). \url{https://arxiv.org/abs/2409.05625v1}.
\bibitem{Hyperbolic} D. Kim, S. H. Lee and S. Lee, Zeta functions of quadratic lattices of a hyperbolic plane, \emph{Math. Nachr.} \textbf{299} (2026), 290--311. \url{https://doi.org/10.1002/mana.70102}.
\bibitem{Maynard} J. Maynard, Primes in arithmetic progressions to large moduli II: well-factorable estimates, \emph{Mem. Amer. Math. Soc.} \textbf{306} (2025), no.~1543. \url{https://doi.org/10.1090/memo/1543}.
\bibitem{Pascadi} A. Pascadi, On the exponents of distribution of primes and smooth numbers, arXiv:2505.00653v2 (2025). \url{https://arxiv.org/abs/2505.00653}.
\bibitem{Rossmann} T. Rossmann, Enumerating submodules invariant under an endomorphism, \emph{Math. Ann.} \textbf{368} (2017), 391--417. \url{https://doi.org/10.1007/s00208-016-1499-6}.
\bibitem{Solomon} L. Solomon, Zeta functions and integral representation theory, \emph{Adv. Math.} \textbf{26} (1977), 306--326. \url{https://doi.org/10.1016/0001-8708(77)90044-5}.
\bibitem{SP} R. Schulze-Pillot, \emph{Lecture Notes on Quadratic Forms and their Arithmetic}, especially \S9.2 (spinor genera and strong approximation), arXiv:2008.12847. \url{https://arxiv.org/abs/2008.12847}.
\bibitem{StevenhagenLenstra} P. Stevenhagen and H. W. Lenstra, Jr., Chebotar\"ev and his density theorem, \emph{Math. Intelligencer} \textbf{18} (1996), no.~2, 26--37. \url{https://doi.org/10.1007/BF03027290}.
\bibitem{Voight} J. Voight, \emph{Quaternion Algebras}, Springer, 2021. \url{https://doi.org/10.1007/978-3-030-56694-4}.
\bibitem{Zeiner} P. Zeiner, Symmetries of coincidence site lattices of cubic lattices, \emph{Z. Kristallogr.} \textbf{220} (2005), 915--925. \url{https://doi.org/10.1524/zkri.2005.220.11_2005.915}.
\bibitem{Zhang} Z. Zhang, Subform zeta functions of positive definite binary quadratic forms of arbitrary conductor, preprint, 2026. \url{https://doi.org/10.2139/ssrn.7390218}.
\end{thebibliography}
\end{document}
